% !TeX root = ../../Book1.tex
\chapter[有界变差函数与有限周长集合]{有界变差函数\\与有限周长集合}
\label{b1:ch:27}
\BookChapterNumber{b1:ch:27}
\begin{chapterabstract}
本章把可积弱梯度推广为有限向量测度，建立有界变差函数的逼近、紧性及余面积公式。
随后从集合示性函数的分布导数定义周长，在约化边界处证明半空间缩放与可整流性，
得到 De Giorgi 结构定理及 Gauss–Green 公式。最后把有界变差函数的 Sobolev 端点估计与集合的等周估计联系起来。
\end{chapterabstract}

\begin{learninggoals}
\begin{enumerate}
\item 从向量测度的欧氏全变差和测试场对偶式识别有界变差函数，区分它与 \(W^{1,1}\) 函数的导数类型。
\item 掌握严格光滑逼近中的单位分解误差，区分局部紧性、全域紧性、弱-*收敛和范数收敛。
\item 从层集截断推导有界变差函数的余面积公式，说明这一推导为何不依赖光滑边界或结构定理。
\item 区分测度论边界、约化边界和任意代表的拓扑边界，从半空间缩放恢复法向及周长面积。
\item 理解函数与集合不等式之间的双向传递，区分本章的有效估计与带最佳常数的尖锐等周定理。
\end{enumerate}
\end{learninggoals}

在一维，阶跃函数虽然没有可积的弱导数，却有一个十分明确的导数：跳跃点上的点质量。
在高维，两个常值区域之间的一道界面也应当贡献变化，而不是因为经典梯度几乎处处为零就被忽略。
有界变差函数空间保留了这类变化，把导数放进有限测度而非仅仅可积函数的范围。
因此它既与前面的 Sobolev 理论连续衔接，又能直接处理不连续函数和集合边界。

本章的证明次序依赖这一选择。先建立变差的分析工具，再用示性函数定义周长；
余面积公式只凭这些工具就能完成。随后再讨论周长究竟集中在哪一种边界上。
“周长等于面积”是这一阶段的结论，不参与此前的定义与逼近。
结构定理中最重要的桥梁是两侧体积下界：方向集中虽然指向一个半空间，
但若没有这项下界，缩放极限仍可能退化为全空间或空集。

Simon 的《Lectures on Geometric Measure Theory》%
\cite[\S\S6 and 14]{Simon1983Lectures}是本章局部变差及边界结构的对读资料。
本书展开其中需要承接的单位分解、平移紧性和图分片步骤；
更系统的有界变差函数理论可参阅 Ambrosio、Fusco 与 Pallara 的%
《Functions of Bounded Variation and Free Discontinuity Problems》%
\cite{Ambrosio2000Functions}。
后者还讨论下一章的跳跃、Cantor 部分与自由间断问题，初读本章不必先掌握这些内容。

% !TeX root = ../../Book1.tex
\section{有界变差函数}
\label{b1:sec:2701}

一维阶跃函数的经典导数几乎处处为零，但它的分布导数在跳跃点留下一个点质量。高维函数也可能把变化集中在超曲面上。若只允许弱导数由 \(L^1\) 函数表示，就会排除这些对象；若允许导数是一份有限 Radon 测度，就能同时保留通常的梯度和集中在薄集上的变化。

\ProseParagraphBreak

本章的函数除非另行说明均为实值，\(\Omega\subseteq\mathbb R^d\) 为开集。导数 \(Du\) 取值于 \(\mathbb R^d\)，其长度由欧氏范数衡量。这里 \(D\) 表示分布导数测度，\(\nabla u\) 则用于确有可积弱梯度的情形；二者不能不加说明地互换。

\ProseParagraphBreak

Simon 的《Lectures on Geometric Measure Theory》\cite[\S6, 6.1]{Simon1983Lectures}从散度测试泛函引入局部有界变差函数。本节先补齐有限维向量测度的变差与极分解，再把这种测试表述同分布导数逐项连接起来。它只调用第6章的 Radon–Nikodym 表示和第11章的 Radon 测度表示，不预借本章后面的几何结构定理。
% 实读本地正式版Simon第6节印刷35--36页，式6.1与局部变差测试表述；
% 该书把向量泛函表示接第4节，本书改经已有标量Riesz--Markov逐分量构造。
% 向量变差的有限分割与连续单位球测试证明在本节完整展开。

\subsection{分布导数为 Radon 测度}
\label{b1:subsec:270101}

先取一个有限向量 Radon 测度 \(\mu=(\mu_1,\ldots,\mu_d)\)，即各分量均为有限实 Radon 符号测度。对 Borel 集 \(A\subseteq\Omega\)，定义其欧氏全变差为
\[
 |\mu|(A)=
 \sup\left\{\sum_{j=1}^N|\mu(A_j)|:
       A=\bigsqcup_{j=1}^N A_j,\ A_j\text{ Borel}\right\}.
\]
这里 \(|\mu(A_j)|\) 是向量的欧氏长度，而不是分量绝对值之和。

\begin{proposition}\label{b1:prop:vector-measure-polar}
令 \(\lambda=\sum_{i=1}^d|\mu_i|\)，并由 Radon–Nikodym 定理写 \(\mu_i=g_i\lambda\)，\(g=(g_1,\ldots,g_d)\)。则
\[
 |\mu|=|g|\lambda.
\]
因此 \(|\mu|\) 是有限正 Radon 测度，且有
\[
 \mu=\sigma_\mu|\mu|,\qquad |\sigma_\mu|=1
 \quad |\mu|\text{-几乎处处}.
\]
这个方向函数在 \(|\mu|\)-几乎处处意义下唯一。上述结论在局部有限情形可逐个内部区域使用，并在重叠处相容。
\end{proposition}
\begin{proof}
对任意有限 Borel 分割，由向量积分的三角不等式，
\[
 \sum_j|\mu(A_j)|=\sum_j\left|\int_{A_j}g\,\mathrm d\lambda\right|
 \le\int_A|g|\,\mathrm d\lambda.
\]
反向取一个有限值向量简单函数 \(s=\sum_j a_j\cdot\mathbf1_{A_j}\)，使
\(\int_A|g-s|\,\mathrm d\lambda<\varepsilon\)。这可以逐分量简单逼近，再取共同的有限分割。于是
\[
 \begin{aligned}
 \sum_j|\mu(A_j)|
 &\ge\sum_j |a_j|\cdot \lambda(A_j)
           -\sum_j\left|\int_{A_j}(g-a_j)\,\mathrm d\lambda\right|\\
 &\ge\int_A|s|\,\mathrm d\lambda-\varepsilon
 \ge\int_A|g|\,\mathrm d\lambda-2\varepsilon.
 \end{aligned}
\]
让误差趋零，得到变差公式，因而也证明了可数可加性。有限 Radon 测度乘以非负可积密度仍为 Radon 测度，这是第11章的密度稳定性。

在 \(g\ne0\) 处取 \(\sigma_\mu=g/|g|\)，其余处取零。测度 \(|\mu|\) 在 \(\{g=0\}\) 上为零，所以所得方向在其几乎处处长度为一，并满足表示式。两个方向若给出同一个向量测度，逐分量用 Radon–Nikodym 唯一性即可证明它们相等。对局部有限分量，把论证限制到内部相对紧开集；唯一性保证局部得到的变差和方向在相交处一致。
\end{proof}

这就是向量测度的极分解，与第6章复测度的方向分解相似。用于控制的 \(\lambda\) 很方便，但真正的欧氏变差是 \(|g|\lambda\)。例如，\(\mu=(m_d,m_d)\) 限制到一个有限体积区域时，其变差是 \(\sqrt2m_d\)，不是 \(2m_d\)。

\begin{definition}\label{b1:def:bv-function}
给定 \(u\in L^1(\Omega)\)。若对每个 \(i=1,\ldots,d\)，分布偏导数都由有限 Radon 符号测度 \(D_i u\) 表示，即
\[
 \int_\Omega u\,\partial_i\varphi\,\mathrm dx
 =-\int_\Omega\varphi\,\mathrm dD_i u
 \quad(\varphi\in C_c^\infty(\Omega)),
\]
则称 \(u\) 为 \(\Omega\) 上的\term[youjie-biancha-hanshu]{有界变差函数}[function of bounded variation]，记 \(u\in BV(\Omega)\)。\foreignidx{bounded variation}\foreignidx{BV}记
\[
 Du\coloneqq(D_1u,\ldots,D_du),
\]
并以 \(|Du|\) 表示其欧氏全变差测度。
\end{definition}
\symidx{\(BV(\Omega)\)：分布导数为有限向量测度的可积函数空间}
\symidx{\(Du\)：有界变差函数的分布导数测度}
\symidx{\(\lvert Du\rvert\)：\(Du\) 的欧氏全变差测度}

若只要求 \(u\in L^1_{\mathrm{loc}}(\Omega)\)，且上述导数测度在每个紧集上具有有限变差，则记 \(u\in BV_{\mathrm{loc}}(\Omega)\)。在内部区域上施加分布测试即可定义这些测度；相邻区域上的表示由测试函数的唯一性相互一致。一个局部有界变差函数不一定在整个 \(\Omega\) 上可积，也不一定有有限的总变差。

\subsection{全变差}
\label{b1:subsec:270102}

变差不仅能由测度分割定义，也能直接从函数通过测试读出。对于任意 \(u\in L^1_{\mathrm{loc}}(\Omega)\) 和开集 \(U\subseteq\Omega\)，规定
\begin{equation}\label{b1:eq:bv-variational-definition}
 V(u,U)\coloneqq\sup\left\{\int_Uu\,\operatorname{div}\varphi\,\mathrm dx:
 \varphi\in C_c^\infty(U;\mathbb R^d),\
 \sup_U|\varphi|\le1\right\}\in[0,\infty].
\end{equation}
每个积分有限，因为测试函数的支集紧含于 \(U\)。把 \(\varphi\) 换成 \(-\varphi\)，说明上确界中也可以对积分取绝对值。

\begin{proposition}\label{b1:prop:bv-variation-duality}
设 \(u\in L^1_{\mathrm{loc}}(\Omega)\)。则 \(V(u,\Omega)<\infty\)，当且仅当分布导数 \(Du\) 是有限向量 Radon 测度。此时对每个开集 \(U\subseteq\Omega\)，
\begin{equation}\label{b1:eq:bv-variation-open-set}
 V(u,U)=|Du|(U).
\end{equation}
若只要求每个 \(\overline U\Subset\Omega\) 的 \(V(u,U)\) 有限，则相应地得到 \(BV_{\mathrm{loc}}\) 的判别和同一局部公式。
\end{proposition}
\begin{proof}
先设 \(V(u,\Omega)<\infty\)。对每个坐标 \(i\)，实线性泛函
\[
 T_i(\psi)=-\int_\Omega u\,\partial_i\psi\,\mathrm dx
 \quad(\psi\in C_c^\infty(\Omega))
\]
满足 \(|T_i(\psi)|\le V(u,\Omega)\|\psi\|_\infty\)，只须在\eqref{b1:eq:bv-variational-definition}中取 \(\varphi=\psi e_i/\|\psi\|_\infty\)。光滑紧支集函数在 \(C_0(\Omega)\) 的上确界范数中稠密：先以紧支集截断，再在离边界有正距离处磨光即可。故 \(T_i\) 唯一延拓为 \(C_0(\Omega)\) 上的有界实泛函。由\cref{b1:thm:riesz-markov-c0}，它由有限 Radon 符号测度 \(D_i u\) 表示，正好满足分布导数关系。

反过来，若 \(Du\) 为有限向量测度，则对任意容许测试场，
\[
 \left|\int_Uu\,\operatorname{div}\varphi\,\mathrm dx\right|
 =\left|\int_U\varphi\cdot\mathrm dDu\right|
 \le\int_U|\varphi|\,\mathrm d|Du|\le|Du|(U).
\]
这先给出 \(V(u,U)\le|Du|(U)\)。

要取得反向不等式，不能直接把可测方向 \(-\sigma_u\) 当作光滑测试场。由有限 Radon 测度的 \(C_c\) 稠密性，逐分量在 \(L^1(|Du|\!\llcorner U)\) 中用连续紧支集向量场逼近 \(-\sigma_u\)。再投影到闭单位球：映射 \(z\mapsto z/\max(1,|z|)\) 连续，且对任意 \(|w|\le1\) 满足
\[
 \left|\frac z{\max(1,|z|)}-w\right|\le2|z-w|.
\]
因此投影后仍能以任意小的积分误差逼近方向，支集仍紧含于 \(U\)，欧氏长度不超过一。将其补零后用非负核小尺度磨光，支集仍在 \(U\) 内，长度由凸性仍不超过一；连续紧支集函数的一致逼近又使积分误差趋零。于是可取光滑场 \(\varphi\)，使
\[
 -\int_U\varphi\cdot\sigma_u\,\mathrm d|Du|
 \ge |Du|(U)-\varepsilon.
\]
由分布恒等式，左端就是 \(\int_Uu\operatorname{div}\varphi\)，从而得到反向不等式。

局部情形在各内部区域实施相同证明。表示测度在重叠处由分布唯一性相等，故可以拼成局部有限的向量 Radon 测度。若开集 \(U\) 上总变差无穷，就先对其中的内部穷尽应用有限情形；下界可以任意大，仍有 \(V(u,U)=|Du|(U)=\infty\)。这也证明了允许无穷值的公式。
\end{proof}

对于 \(u\in W^{1,1}(\Omega)\)，各测度 \(D_i u\) 都有密度 \(\partial_i u\)，前面的向量分解给出
\begin{equation}\label{b1:eq:sobolev-bv-variation}
 Du=\nabla u\,m_d,\qquad
 |Du|=|\nabla u|\,m_d.
\end{equation}
因此 \(W^{1,1}(\Omega)\subseteq BV(\Omega)\)，且光滑函数的变差就是梯度长度的积分。一般有界变差函数的导数测度可能含有相对于体积的奇异部分，不能把\eqref{b1:eq:sobolev-bv-variation}不加条件地用于所有有界变差函数。

\ProseParagraphBreak

这一测试表述也是下半连续性的来源：固定测试场时，积分只依赖 \(u\) 在一个内部紧集上的 \(L^1\) 值；取所有这类连续表达式的上确界，便得到可通过极限保留的下界。下一节将把这个事实连同光滑逼近一起用于紧性。

\subsection{有界变差函数空间的范数}
\label{b1:subsec:270103}

在 \(BV(\Omega)\) 上规定
\[
 \|u\|_{BV(\Omega)}\coloneqq\|u\|_{L^1(\Omega)}+|Du|(\Omega).
\]
向量测度的三角不等式保证它满足范数三角不等式，零阶项则保证范数为零只对应零等价类。

\begin{proposition}\label{b1:prop:bv-complete}
上述范数使 \(BV(\Omega)\) 成为 Banach 空间。
\end{proposition}
\begin{proof}
设 \((u_n)\) 是 Cauchy 列。由 \(L^1\) 完备性，存在 \(u\in L^1(\Omega)\)，使 \(u_n\to u\) 于 \(L^1\)。同时
\[
 |D_i u_n-D_i u_m|(\Omega)
 \le|D(u_n-u_m)|(\Omega)\longrightarrow0.
\]
第11章的有限 Radon 测度空间在全变差范数下完备，因此每个分量都趋于某个有限 Radon 符号测度 \(\mu_i\)。对任意 \(\varphi\in C_c^\infty(\Omega)\)，
\[
 \int_\Omega u\,\partial_i\varphi
 =\lim_n\int_\Omega u_n\partial_i\varphi
 =-\lim_n\int_\Omega\varphi\,\mathrm dD_i u_n
 =-\int_\Omega\varphi\,\mathrm d\mu_i.
\]
故 \(D_i u=\mu_i\)，而 \(u\in BV(\Omega)\)。再由
\[
 |D(u_n-u)|(\Omega)
 \le\sum_{i=1}^d|D_i u_n-\mu_i|(\Omega)\longrightarrow0,
\]
得到原范数中的收敛。
\end{proof}

局部化仍然可用，但乘法公式中需要同时保留测度项与体积项。

\begin{proposition}\label{b1:prop:bv-product-rule}
设 \(u\in BV(\Omega)\)。若 \(v\) 在 \(\Omega\) 中局部 Lipschitz，并满足 \(v,\nabla v\in L^\infty(\Omega)\)，则
\[
 v\cdot u\in BV(\Omega),\qquad
 D(v\cdot u)=v\cdot Du+u\cdot\nabla v\,m_d,
\]
且
\[
 |D(vu)|(\Omega)
 \le\|v\|_\infty|Du|(\Omega)
       +\|\nabla v\|_\infty\|u\|_1.
\]
第一项使用 \(v\) 的连续点值，第二项使用其几乎处处弱梯度。相应局部结论也成立。
\end{proposition}
\begin{proof}
先取光滑 \(v\)。对任意光滑紧支集测试函数，将乘积 \(v\cdot\varphi\) 代入 \(Du\) 的定义，再用普通乘积求导：
\[
 \int_\Omega u\cdot v\,\partial_i\varphi
 =\int_\Omega u\,\partial_i(v\cdot\varphi)
       -\int_\Omega u\cdot\varphi\,\partial_i v
 =-\int_\Omega v\cdot\varphi\,\mathrm dD_i u
       -\int_\Omega u\cdot\varphi\,\partial_i v.
\]
这就是所需测度公式。若 \(v\) 只局部 Lipschitz，固定测试函数的支集，在其邻域中磨光 \(v\)。磨光函数局部一致趋于 \(v\)，梯度局部一致有界，且沿一个子列几乎处处趋于 \(\nabla v\)。因此测度积分由局部一致收敛通过极限，体积积分由 \(|u|\) 控制收敛通过极限。每个测试函数都得到相同公式，故分布表示成立。

最后，对 Borel 集上的向量积分取变差，得到
\[
 |D(vu)|\le |v|\,|Du|+|u|\cdot|\nabla v|\,m_d.
\]
积分即得估计，且零阶积 \(v\cdot u\) 可积。局部版本限制到内部区域即可。
\end{proof}

特别地，乘以内部光滑紧支集截断后，可以在域外补零，得到全空间有界变差函数。证明只需在截断支集附近代入测试函数，因而不会产生边界项。这个内部补零与在任意边界直接补零不同；后者可能增加一份边界测度，甚至使总变差失去控制。

\subsection{基本例子}
\label{b1:subsec:270104}

在区间 \((-1,1)\) 上取 \(u=\mathbf1_{(0,1)}\)。对任意测试函数，
\[
 \int_{-1}^1u\varphi'\,\mathrm dx
 =\int_0^1\varphi'\,\mathrm dx=-\varphi(0).
\]
因此 \(Du=\delta_0\)，\(|Du|((-1,1))=1\)。这个函数属于 \(BV\)，却不属于 \(W^{1,1}\)，因为点质量没有 \(L^1\) 密度。导数测度记录了一次跳跃的大小，而不是把它藏进“经典导数几乎处处为零”的表述中。

\ProseParagraphBreak

高维的长方体可用相同计算处理。令 \(Q=\prod_{i=1}^d(a_i,b_i)\)。在全空间中，
\[
 \int_Q\partial_i\varphi\,\mathrm dx
 =\int_{\prod_{j\ne i}(a_j,b_j)}
       \bigl(\varphi(x',b_i)-\varphi(x',a_i)\bigr)\,\mathrm dx'.
\]
所以 \(D_i\mathbf1_Q\) 在下侧面取正的面测度，在上侧面取负的面测度。除边棱这个面测度零集以外，各侧面互不重叠；导数方向在每个面上均为单位法向的负号。因此
\[
 |D\mathbf1_Q|(\mathbb R^d)
 =2\sum_{i=1}^d\prod_{j\ne i}(b_j-a_j).
\]
这是边界各面的面积之和。当 \(d=1\) 时空乘积取一，恢复一个有界区间两端的总变差二。到\secref{b1:sec:2704}，我们会把这种现象扩展为一般集合的周长。

\ProseParagraphBreak

变化也可能既不由体积密度给出，又不集中在离散跳点。令 \(\mu_C\) 为标准 Cantor 概率测度，\(c(x)=\mu_C([0,x])\)，\(0<x<1\)。对 \(\varphi\in C_c^\infty(0,1)\)，Fubini 定理给出
\[
 \begin{aligned}
 \int_0^1c(x)\varphi'(x)\,\mathrm dx
 &=\int_{[0,1]}\left(\int_t^1\varphi'(x)\,\mathrm dx\right)
                                      \mathrm d\mu_C(t)\\
 &=-\int_{[0,1]}\varphi(t)\,\mathrm d\mu_C(t).
 \end{aligned}
\]
端点不带原子，故在开区间内有 \(Dc=\mu_C\)，总变差为一。函数 \(c\) 连续，在 Cantor 集的补区间上局部常数，所以经典导数几乎处处为零；但导数测度仍保留了非零的奇异连续变化。这是后面 Cantor 部分的最基本模型。

\ProseParagraphBreak

这些例子也说明，有界变差函数空间由积分与分布刻画，不是对任意点态代表计算振幅总和。把零函数在有理点上改为一，不改变其 \(L^1\) 类和分布导数，所以在多维定义下仍代表 \(BV\) 空间的零元。若回到一维的点态分割变差，则需先选择相应的规范代表；不能把任意零测改动后的点态分割和与这里的 \(|Du|\) 混为一谈。

% !TeX root = ../../Book1.tex
\section{有界变差函数的逼近与紧性}
\label{b1:sec:2702}

对 Sobolev 函数，磨光可以同时逼近函数及其弱梯度；对一般有界变差函数，导数可能含有集中在零测集上的部分，光滑梯度不能按总变差范数逼近这种测度。本节保留函数的 \(L^1\) 收敛和导数的总量收敛，并用导数测度控制平移，得到局部紧性。所有函数均取实值，\(|Du|\) 使用前一节的欧氏全变差。

Simon 的《Lectures on Geometric Measure Theory》\cite[\S6, 6.2 LEMMA and 6.3 THEOREM]{Simon1983Lectures}以磨光及变差对偶式组织有界变差函数的逼近与紧性。下面的全域严格逼近另外展开单位分解的误差控制；局部紧性的抽选步骤则调用已经证明的 \(L^1\) 平移判据。

% 实读 Simon1983Lectures 正式公开本§6，印刷35--38页：
% 6.1局部BV的测度表示、6.2局部磨光及变差测度收敛、6.3局部紧性；
% 核看第37页。原书6.3将平滑函数的紧性细节交给GT，本节改用已证FK。
% 任意开集的全域严格逼近按22.1局部有限分区路线独立展开，
% 特别控制rho_i*(u grad psi_i)-u grad psi_i，不声称直接来自Simon完整证明。
% 不使用一般BV边界延拓；补零只对内部光滑截断或另列已知BV补零假设。

\subsection{光滑逼近}
\label{b1:subsec:270201}

\begin{definition}\label{b1:def:bv-strict-convergence}
设 \(u_j,u\in BV(\Omega)\)。若
\[
 \|u_j-u\|_{L^1(\Omega)}\longrightarrow0,\quad
 |Du_j|(\Omega)\longrightarrow|Du|(\Omega),
\]
则称 \(u_j\) 严格收敛到 \(u\)；这种收敛称为\term[youjie-biancha-kongjian-yange-shoulian]{有界变差空间中的严格收敛}[strict convergence in BV]。
\end{definition}

这一定义比较的是两个导数测度的总变差数值，并没有要求
\(|D(u_j-u)|(\Omega)\to0\)。先记录内部截断及全空间磨光的公式。

\begin{lemma}\label{b1:lem:bv-cutoff-and-mollification}
设 \(u\in BV_{\mathrm{loc}}(\Omega)\)、\(\eta\in C_c^\infty(\Omega)\)，把 \(\eta u\) 在 \(\Omega\) 外补零，记为 \(v\)。则 \(v\in BV(\mathbb R^d)\)，且
\begin{equation}\label{b1:eq:bv-cutoff-derivative}
 Dv=\eta Du+u\nabla\eta\,m_d,
 \quad
 |Dv|(\mathbb R^d)
 \leq\int_\Omega|\eta|\,\mathrm d|Du|
       +\int_\Omega|u|\cdot|\nabla\eta|\,\mathrm dx,
\end{equation}
右侧两项测度均在域外补零。

若 \(v\in BV(\mathbb R^d)\)，用非负、偶、单位积分的磨光函数定义
\(v_\varepsilon=\rho_\varepsilon*v\)，则
\[
 v_\varepsilon\in C^\infty(\mathbb R^d)\cap W^{1,1}(\mathbb R^d),
 \quad v_\varepsilon\longrightarrow v\ \text{于 }L^1,
\]
并有
\begin{equation}\label{b1:eq:bv-convolution-gradient}
 \nabla v_\varepsilon=\rho_\varepsilon*Dv,\quad
 \int_{\mathbb R^d}|\nabla v_\varepsilon|\,\mathrm dx
 \leq|Dv|(\mathbb R^d).
\end{equation}
\end{lemma}
\begin{proof}
对任意 \(\Phi\in C_c^\infty(\mathbb R^d;\mathbb R^d)\)，函数
\(\eta\Phi\) 的支集紧含于 \(\Omega\)。乘积公式给出
\[
 \begin{aligned}
 \int_{\mathbb R^d}v\,\operatorname{div}\Phi
 &=\int_\Omega u\,\operatorname{div}(\eta\Phi)
       -\int_\Omega u\nabla\eta\cdot\Phi\\
 &=-\int_\Omega\eta\Phi\cdot\mathrm dDu
       -\int_\Omega u\nabla\eta\cdot\Phi.
 \end{aligned}
\]
两项导数测度均有限，得到补零后的导数公式；总变差的三角不等式给出估计。这里没有跨越 \(\partial\Omega\) 的边界项，因为 \(\eta\) 的支集留在域内。

对任意有限向量测度 \(\mu\)，定义
\[
 (\rho_\varepsilon*\mu)(x)
 =\int_{\mathbb R^d}\rho_\varepsilon(x-y)\,\mathrm d\mu(y).
\]
由向量测度的极分解和积分三角不等式，
\[
 \bigl|(\rho_\varepsilon*\mu)(x)\bigr|
 \leq\int\rho_\varepsilon(x-y)\,\mathrm d|\mu|(y),
 \quad
 \int|\rho_\varepsilon*\mu|\,\mathrm dx\leq|\mu|(\mathbb R^d).
\]
后一式由 Tonelli 定理及核的单位积分得到。

函数 \(v_\varepsilon\) 的光滑性由卷积求导给出，而 \(L^1\) 收敛来自\cref{b1:prop:approximate-identity-lp}。利用核的偶性及 Fubini 定理，对测试向量场 \(\Phi\) 有
\[
 \begin{aligned}
 \int v_\varepsilon\operatorname{div}\Phi
 &=\int v\,\operatorname{div}(\rho_\varepsilon*\Phi)\\
 &=-\int(\rho_\varepsilon*\Phi)\cdot\mathrm dDv
 =-\int\Phi(x)\cdot(\rho_\varepsilon*Dv)(x)\,\mathrm dx.
 \end{aligned}
\]
因此经典梯度为 \(\rho_\varepsilon*Dv\)。刚证的向量测度卷积估计说明它可积，并给出\eqref{b1:eq:bv-convolution-gradient}。
\end{proof}

\begin{theorem}\label{b1:thm:bv-strict-approximation}
设 \(\Omega\subseteq\mathbb R^d\) 为任意非空开集，\(u\in BV(\Omega)\)。存在
\[
 u_j\in C^\infty(\Omega)\cap W^{1,1}(\Omega)
\]
严格收敛到 \(u\)，也就是
\[
 u_j\to u\ \text{于 }L^1(\Omega),\quad
 \int_\Omega|\nabla u_j|\,\mathrm dx\longrightarrow|Du|(\Omega).
\]
这里不要求 \(u_j\) 紧支集，也不要求边界正则。
\end{theorem}
\begin{proof}
给定 \(\varepsilon>0\)。沿用\cref{b1:thm:meyers-serrin}证明中的局部有限开层及光滑单位分解，取
\[
 \psi_i\geq0,\quad \sum_i\psi_i=1,\quad
 K_i=\operatorname{supp}\psi_i\Subset V_{j(i)},\quad
 \overline{V_{j(i)}}\Subset\Omega.
\]
每个紧集只遇到有限个开层 \(V_j\)，每一层只分配有限多个 \(\psi_i\)。

令 \(f_i=\psi_i u\)，并补零到全空间。由前一引理，
\[
 Df_i=\mu_i+g_i\,m_d,\quad
 \mu_i=\psi_iDu,\quad g_i=u\nabla\psi_i.
\]
每个 \(f_i,g_i\) 都属于全空间 \(L^1\)，且支集包含于 \(K_i\)。选择足够小的半径 \(\delta_i>0\)，同时满足
\[
 K_i+\overline B(0,\delta_i)\subset V_{j(i)},\quad
 \|\rho_{\delta_i}*f_i-f_i\|_1<\varepsilon2^{-i-1},\quad
 \|\rho_{\delta_i}*g_i-g_i\|_1<\varepsilon2^{-i-1}.
\]
支集条件由紧集到补集的正距离保证；两个误差条件由 \(L^1\) 磨光逼近保证，因而可以通过进一步缩小同一个半径同时实现。

定义
\[
 v=\sum_i\rho_{\delta_i}*f_i.
\]
每个紧集只遇到有限个磨光后的支集，所以此和在局部为有限个光滑函数之和，故 \(v\in C^\infty(\Omega)\)。又因
\[
 \sum_i\|f_i\|_1=\int_\Omega|u|\sum_i\psi_i=\|u\|_1,
\]
卷积后的级数也在 \(L^1\) 中绝对收敛，并且
\[
 \|v-u\|_1
 \leq\sum_i\|\rho_{\delta_i}*f_i-f_i\|_1<\varepsilon/2.
\]

梯度需要另作估计。在任意内部紧集上，单位分解满足
\(\sum_i\nabla\psi_i=0\)。因此局部逐项求导给出
\begin{equation}\label{b1:eq:bv-partition-commutator}
 \nabla v
 =\sum_i\rho_{\delta_i}*\mu_i
       +\sum_i(\rho_{\delta_i}*g_i-g_i).
\end{equation}
第一组项的全域 \(L^1\) 范数可求和，因为非负权重给出
\[
 \sum_i\int_{\mathbb R^d}|\rho_{\delta_i}*\mu_i|
 \leq\sum_i|\mu_i|(\mathbb R^d)
 =\sum_i\int_\Omega\psi_i\,\mathrm d|Du|
 =|Du|(\Omega).
\]
第二组误差的 \(L^1\) 范数之和小于 \(\varepsilon/2\)。所以
\[
 v\in W^{1,1}(\Omega),\quad
 \int_\Omega|\nabla v|\leq|Du|(\Omega)+\varepsilon/2.
\]
这里没有要求 \(\sum_i\|u\nabla\psi_i\|_1<\infty\)；相消先在局部进行，全域求和只作用于已经控制的误差。

令 \(\varepsilon\downarrow0\)，得到所需的 \(L^1\) 逼近列及梯度积分的上极限估计。为取得反向不等式，对每个
\(\Phi\in C_c^\infty(\Omega;\mathbb R^d)\)、\(|\Phi|\leq1\)，有
\[
 \int_\Omega u\,\operatorname{div}\Phi
 =\lim_j\int_\Omega u_j\operatorname{div}\Phi
 \leq\liminf_j\int_\Omega|\nabla u_j|.
\]
对 \(\Phi\) 取上确界，由变差对偶式得到
\(|Du|(\Omega)\leq\liminf_j\int_\Omega|\nabla u_j|\)，完成证明。
\end{proof}

有界变差空间中的严格收敛确实弱于 \(BV\) 范数收敛。例如在 \((-1,1)\) 上取
\(u=\mathbf1_{(0,1)}\)，则 \(Du=\delta_0\)。令
\[
 u_\varepsilon(x)=\int_{-\infty}^x\rho_\varepsilon(t)\,\mathrm dt,
 \quad 0<\varepsilon<1.
\]
这些函数光滑，与 \(u\) 只在 \((-\varepsilon,\varepsilon)\) 内不同，且
\(\int_{-1}^1|u_\varepsilon'|=1=|Du|\bigl((-1,1)\bigr)\)。因此它们严格收敛，但
\(\rho_\varepsilon\,m_1\) 与 \(\delta_0\) 相互奇异，所以
\[
 |D(u_\varepsilon-u)|\bigl((-1,1)\bigr)=2.
\]
不能把光滑严格逼近改称一般有界变差函数的光滑范数稠密性。

\subsection{下半连续性}
\label{b1:subsec:270202}

前面最后一步只使用了紧支集测试，因此它实际上是一个局部结论。

\begin{theorem}[变差的下半连续性]\label{b1:thm:bv-lower-semicontinuity}
设 \(u_j,u\in L^1_{\mathrm{loc}}(\Omega)\)，且
\(u_j\to u\) 于 \(L^1_{\mathrm{loc}}(\Omega)\)。对任意开集 \(U\subseteq\Omega\)，有
\[
 V(u,U)\leq\liminf_{j\to\infty}V(u_j,U),
\]
其中 \(V\) 是\cref{b1:prop:bv-variation-duality}中的扩展变差，允许取无穷。

若右端在每个内部相对紧开集上有限，则 \(u\in BV_{\mathrm{loc}}(\Omega)\)。在导数测度存在时，上式就是
\[
 |Du|(U)\leq\liminf_j|Du_j|(U).
\]
若另有 \(u\in L^1(\Omega)\) 及 \(\liminf_jV(u_j,\Omega)<\infty\)，则 \(u\in BV(\Omega)\)。
\end{theorem}
\begin{proof}
任取 \(\Phi\in C_c^\infty(U;\mathbb R^d)\)、\(|\Phi|\leq1\)。因为
\(\operatorname{div}\Phi\) 有界且紧支集，局部 \(L^1\) 收敛给出
\[
 \int_Uu\,\operatorname{div}\Phi
 =\lim_j\int_Uu_j\operatorname{div}\Phi
 \leq\liminf_jV(u_j,U).
\]
右端不依赖 \(\Phi\)，故对测试场取上确界，得到第一项。其余断言由局部与全局的有限变差判别及 \(V=|Du|\) 得到。
\end{proof}

对层集示性函数也可以直接用这个版本：即使 \(\{u>t\}\) 的体积无穷，
\(\mathbf1_{\{u>t\}}\) 仍局部可积，扩展变差依然有定义。全域
\(L^1\) 是进入 \(BV(\Omega)\) 的额外要求，不是定义局部周长的必要条件。

\subsection{有界变差函数的紧性定理}
\label{b1:subsec:270203}

导数是测度时，平移仍能由它的总变差控制。

\begin{lemma}\label{b1:lem:bv-translation-estimate}
若 \(v\in BV(\mathbb R^d)\)，则对每个 \(h\in\mathbb R^d\)，有
\[
 \|\tau_hv-v\|_{L^1(\mathbb R^d)}
 \leq |h|\cdot|Dv|(\mathbb R^d),
 \quad \tau_hv(x)=v(x-h).
\]
\end{lemma}
\begin{proof}
先对全空间磨光 \(v_\varepsilon\) 使用沿线段的微积分基本定理：
\[
 |v_\varepsilon(x-h)-v_\varepsilon(x)|
 \leq|h|\int_0^1|\nabla v_\varepsilon(x-th)|\,\mathrm dt.
\]
积分换序及平移不变性给出
\[
 \|\tau_hv_\varepsilon-v_\varepsilon\|_1
 \leq|h|\int|\nabla v_\varepsilon|
 \leq|h|\cdot|Dv|(\mathbb R^d).
\]
由 \(v_\varepsilon\to v\) 于 \(L^1\)，以及平移的等距性，左端在
\(\varepsilon\downarrow0\) 时趋于 \(\|\tau_hv-v\|_1\)，从而完成证明。
\end{proof}

\begin{theorem}[有界变差函数的局部紧性]\label{b1:thm:bv-local-compactness}
设 \(u_j\in BV_{\mathrm{loc}}(\Omega)\)，且对每个
\(\overline U\Subset\Omega\) 的开集 \(U\)，都有
\[
 \sup_j\bigl(\|u_j\|_{L^1(U)}+|Du_j|(U)\bigr)<\infty.
\]
则存在子列及 \(u\in BV_{\mathrm{loc}}(\Omega)\)，使
\[
 u_{j_k}\longrightarrow u\quad\text{于 }L^1_{\mathrm{loc}}(\Omega).
\]
对任意开集 \(U\subseteq\Omega\)，所得极限满足
\(|Du|(U)\leq\liminf_k|Du_{j_k}|(U)\)。
\end{theorem}
\begin{proof}
取递增开集穷竭 \(U_m\)，使
\(\overline{U_m}\Subset U_{m+1}\)，并取
\(\eta_m\in C_c^\infty(U_{m+1})\)、\(0\leq\eta_m\leq1\)，在 \(U_m\) 上等于 \(1\)。把 \(\eta_mu_j\) 补零到全空间，记为 \(v_{m,j}\)。由截断引理及假设，
\[
 \sup_j\bigl(\|v_{m,j}\|_1+|Dv_{m,j}|(\mathbb R^d)\bigr)<\infty
\]
对每个固定 \(m\) 成立，且这组函数都在
\(\operatorname{supp}\eta_m\) 外为零。

共同紧支集保证尾部条件，前一引理保证平移误差一致趋零。由\cref{b1:thm:frechet-kolmogorov-lp}取 \(p=1\)，对固定 \(m\)，族
\((v_{m,j})_j\) 在全空间 \(L^1\) 中相对紧，因而 \((u_j)_j\) 在
\(L^1(U_m)\) 中有收敛子列。

逐次为 \(m=1,2,\ldots\) 选取子列，再取对角子列。其在每个 \(U_m\) 上都有 \(L^1\) 极限，不同区域的极限在交集上几乎处处相同，因为
\(L^1\) 极限唯一。因此这些极限定义了一个 \(u\in L^1_{\mathrm{loc}}(\Omega)\)。局部变差界及下半连续性保证 \(u\in BV_{\mathrm{loc}}\)，并给出全部开集上的所述不等式。
\end{proof}

若原序列的 \(BV(\Omega)\) 范数一致有界，局部紧性得到的极限也属于全域
\(BV(\Omega)\)：可以再抽出几乎处处收敛子列，用 Fatou 引理控制其
\(L^1\) 范数，再用变差下半连续性控制 \(|Du|(\Omega)\)。但这个结论尚未给出全域 \(L^1\) 收敛，质量仍可能移到穷竭所遗漏的区域。

\begin{corollary}\label{b1:cor:bv-global-compactness-with-tails}
设 \((u_j)\) 在 \(BV(\Omega)\) 中一致有界。若存在递增紧集
\(K_m\Subset\Omega\)，使
\[
 \lim_{m\to\infty}\sup_j
       \int_{\Omega\setminus K_m}|u_j|\,\mathrm dx=0,
\]
则它有全域 \(L^1(\Omega)\) 收敛子列。

特别地，若 \(\Omega\) 有界，只需另有边界附近的统一质量控制
\[
 \lim_{\delta\downarrow0}\sup_j
 \int_{\{\,x\in\Omega:\operatorname{dist}(x,\Omega^c)<\delta\,\}}
       |u_j(x)|\,\mathrm dx=0.
\]
另一项充分条件是：\(\Omega\) 有界，而且补零函数 \(E_0u_j\) 属于
\(BV(\mathbb R^d)\)，并在该空间中一致有界。
\end{corollary}
\begin{proof}
先取局部 \(L^1\) 收敛子列，再抽出几乎处处收敛子列，极限记为
\(u\in BV(\Omega)\)。Fatou 引理使极限继承同一个尾部界。给定
\(\varepsilon>0\)，取 \(m\)，使序列及极限在 \(\Omega\setminus K_m\) 上的 \(L^1\) 范数都小于 \(\varepsilon\)。于是
\[
 \|u_j-u\|_{L^1(\Omega)}
 \leq\|u_j-u\|_{L^1(K_m)}+2\varepsilon.
\]
第一项由局部收敛趋于零，再令 \(\varepsilon\downarrow0\)。

当 \(\Omega\) 有界时，可取
\(K_m=\{\,x\in\Omega:\operatorname{dist}(x,\Omega^c)\geq1/m\,\}\)；
这些集合紧含于 \(\Omega\)，并穷竭它，所以边界条件推出前述尾部条件。

最后一种情形直接在全空间应用平移引理和 Fréchet--Kolmogorov 定理：补零函数具有共同有界支承区域，且全空间导数总变差一致有界。所得全空间 \(L^1\) 收敛限制到 \(\Omega\) 即可。
\end{proof}

仅有“\(\Omega\) 有界”不能代替这些条件。第\ref{b1:ex:25-infinitely-many-components}题中，取 \(p=1\)，在越来越小的不同连通分量上放置归一化常数，得到
\[
 \|u_j\|_{L^1(\Omega)}=1,\quad Du_j=0,\quad
 \|u_j-u_k\|_{L^1(\Omega)}=2\quad(j\ne k).
\]
这同时是有界变差函数空间中的有界列，却没有全域 \(L^1\) 收敛子列。对有界
Lipschitz 区域，若要通过延拓证明无附加条件的有界变差函数的全局紧性，还须先建立有界变差函数的延拓定理；第22章的 \(W^{1,1}\) 延拓算子并未定义在一般有界变差函数上，不能直接套用。

\subsection{弱-*收敛}
\label{b1:subsec:270204}

局部 \(L^1\) 收敛首先决定分布导数的极限，变差界再把光滑测试扩展为连续测试。

\begin{proposition}\label{b1:prop:bv-distributional-weak-star}
设 \(u_j,u\in BV_{\mathrm{loc}}(\Omega)\)，\(u_j\to u\) 于
\(L^1_{\mathrm{loc}}(\Omega)\)，且导数变差在每个内部紧集的邻域上一致有界。则对每个
\(\Phi\in C_c(\Omega;\mathbb R^d)\)，有
\[
 \int_\Omega\Phi\cdot\mathrm dDu_j
 \longrightarrow\int_\Omega\Phi\cdot\mathrm dDu.
\]
若进一步有 \(\sup_j|Du_j|(\Omega)<\infty\)，则 \(|Du|(\Omega)<\infty\)，并且测试函数可以取为任意 \(\Phi\in C_0(\Omega;\mathbb R^d)\)。
\end{proposition}
\begin{proof}
若 \(\Phi\) 光滑且紧支集，由分布导数定义，
\[
 \int\Phi\cdot\mathrm dDu_j
 =-\int u_j\operatorname{div}\Phi
 \longrightarrow-\int u\operatorname{div}\Phi
 =\int\Phi\cdot\mathrm dDu.
\]
给定连续紧支集 \(\Phi\)，可把它补零后磨光，得到统一收敛于它的光滑向量场 \(\Phi_\ell\)，并使这些场的支集留在某个固定
\(K\Subset\Omega\) 内。若 \(M_K=\sup_j|Du_j|(K)<\infty\)，则
\[
 \left|\int(\Phi-\Phi_\ell)\cdot\mathrm dDu_j\right|
 \leq M_K\|\Phi-\Phi_\ell\|_\infty.
\]
极限测度有同样的估计。先固定 \(\ell\) 令 \(j\to\infty\)，再令
\(\ell\to\infty\)，得到连续紧支集测试的结论。

若总变差一致有界，下半连续性给出 \(|Du|(\Omega)<\infty\)。由\cref{b1:prop:cc-dense-c0}，连续紧支集函数在 \(C_0\) 中一致稠密；用全域总变差界重复刚才的误差估计，就得到全部 \(C_0\) 测试。
\end{proof}

\begin{definition}\label{b1:def:bv-weak-star-convergence}
设 \(u_j,u\in BV(\Omega)\)。若 \(u_j\to u\) 于 \(L^1(\Omega)\)，且有限向量测度 \(Du_j\) 对全部 \(C_0(\Omega;\mathbb R^d)\) 测试弱-*收敛到 \(Du\)，则称这种收敛为\term[youjie-biancha-kongjian-ruoxing-shoulian]{有界变差空间中的弱-*收敛}[weak-* convergence in BV]。
\end{definition}

因此，对于有界变差函数空间中的有界列，一旦已经取得全域 \(L^1\) 收敛，导数测度的弱-*收敛就随之成立。只取得局部 \(L^1\) 收敛时，仍有前一命题的测度测试结论，却不能省去本定义中的全域 \(L^1\) 条件。

弱-*收敛不保证总变差数值收敛。边界逃逸给出一个例子：在
\(\Omega=(0,1)\) 上，令 \(u_j=\mathbf1_{(0,1/j)}\)。则
\[
 u_j\to0\ \text{于 }L^1,\quad
 Du_j=-\delta_{1/j}\overset{*}{\rightharpoonup}0,\quad
 |Du_j|(\Omega)=1.
\]
这里 \(C_0\bigl((0,1)\bigr)\) 测试在边界附近趋零；常数函数 \(1\) 不属于这个测试空间。若改用全部 \(C_b\) 测试，就换成了\chapref{b1:ch:12}的另一种测度收敛。

振荡也会造成变差损失，即使不把质量集中到边界。取
\(\Omega=(0,2\pi)\)、\(u_j(x)=j^{-1}\sin(jx)\)。函数在 \(L^1\) 中趋零，导数为 \(\cos(jx)\,m_1\)，故由前一命题弱-*趋零，但
\[
 |Du_j|(\Omega)=\int_0^{2\pi}|\cos(jx)|\,\mathrm dx=4.
\]
这些变差测度在边界宽度为 \(\delta\) 的区域内质量至多 \(2\delta\)，所以即使控制了边界质量，也不能从导数的弱-*收敛推出总变差收敛。严格逼近定理额外取得的，正是后一个数值极限。

% !TeX root = ../../Book1.tex
\section{BV 的余面积公式}
\label{b1:sec:2703}

一个实值函数可以由全部超水平集重建。对于 Sobolev 函数，第%
\ref{b1:sec:1802}节用梯度长度补偿水平集上的面积；对于 BV 函数，
梯度不再一定是函数，但超水平集的边界仍可由示性函数的分布导数描述。
本节证明：函数的总变差恰好是各层周长的积分。证明先用窄条带截断处理
\(W^{1,1}\) 函数，再借严格逼近通过极限，不以有限周长集合的深层边界结构为前提。

% 实际来源范围：本地 Simon1983Lectures 正式扫描本，第6节印刷34--40页，
% 6.1的局部变差对偶、6.2的磨光及6.4--6.6的Poincare型估计。
% 该范围没有提供本节BV余面积的完整证明；下面条带截断、可数测试族和
% 有符号分层的证明由本书已证工具独立组织，不冒称来自原书某一未读定理。
% 依赖：27.1变差对偶；27.2扩展变差L1loc下半连续与任意开集严格逼近；
% 21章Sobolev截断。第18章仅作几何对照，不作为本证明的必要黑箱。

\subsection{超水平集}
\label{b1:subsec:270301}

以下固定开集 \(\Omega\subseteq\mathbb R^d\)。为实值可测函数 \(u\) 记
\[
 E_t=\{\,x\in\Omega:u(x)>t\,\},\quad t\in\mathbb R.
\]
改动 \(u\) 的一个 Lebesgue 零测集，只会在同一个零测集内改动每个 \(E_t\)。
示性函数的分布导数不受这种改动影响，所以后面的周长由原来的几乎处处类决定。
需要联合可测性时，可以先选一个 Borel 代表。

对有限实数 \(a,b\)，有两个基本恒等式：
\begin{equation}\label{b1:eq:bv-signed-layer-cake}
 a=\int_{\mathbb R}
       \bigl(\mathbf1_{\{a>t\}}-\mathbf1_{\{0>t\}}\bigr)\,\mathrm dt,
 \quad
 |a-b|=\int_{\mathbb R}
       \bigl|\mathbf1_{\{a>t\}}-\mathbf1_{\{b>t\}}\bigr|\,\mathrm dt.
\end{equation}
第一式在 \(a\ge0\) 时积分区间为 \((0,a)\)，在 \(a<0\) 时则在
\((a,0)\) 上积分 \(-1\)；第二式的被积函数只在两个端点之间等于 \(1\)。
负值情形必须保留第一式中的基准项，不能将有符号函数写成所有超水平集示性函数的直接积分。

由第二式及 Tonelli 定理，对 \(u,v\in L^1(U)\)，有
\begin{equation}\label{b1:eq:bv-level-set-l1-distance}
 \int_{\mathbb R}
 m_d\bigl(\{u>t\}\mathbin{\triangle}\{v>t\}\bigr)\,\mathrm dt
   =\|u-v\|_{L^1(U)},
\end{equation}
其中所有层集均取在 \(U\) 内。若 \(u_j\to u\) 于 \(L^1(U)\)，
抽取满足 \(\sum_j\|u_j-u\|_1<\infty\) 的子列，就得到对几乎每个 \(t\)，
\[
 \mathbf1_{\{u_j>t\}}\longrightarrow\mathbf1_{\{u>t\}}
       \quad\text{于 }L^1(U).
\]
这里的子列对几乎所有层值共同有效；它不是为每个 \(t\) 分别抽取的子列。

\subsection{周长}
\label{b1:subsec:270302}

在任意开集 \(U\subseteq\Omega\) 上，沿用前节的扩展变差记号
\[
 V(f,U)=\sup\left\{\,
     \int_U f\,\operatorname{div}\varphi\,\mathrm dx:
     \varphi\in C_c^1(U;\mathbb R^d),\
     \|\varphi\|_\infty\le1\,\right\}.
\]
\(C_c^1\) 测试场可在共同紧集内磨光，并在必要时按上确界范数归一化；
因此这个上确界与前节使用 \(C_c^\infty\) 的定义相同。
这个数允许为无穷；\(f\) 只需局部可积。测试场的大小始终用欧氏范数。

\begin{definition}\label{b1:def:perimeter}
设 \(E\subseteq\Omega\) 为 Lebesgue 可测集。定义它在开集 \(U\subseteq\Omega\) 内的%
\term[zhouchang]{周长}[perimeter]为
\[
 P(E,U)=V(\mathbf1_E,U).
\]
若 \(P(E,\Omega)<\infty\)，则称 \(E\) 为 \(\Omega\) 内的%
\term[youxian-zhouchang-ji]{有限周长集}[set of finite perimeter]。
若对每个满足 \(\overline U\Subset\Omega\) 的开集 \(U\) 都有 \(P(E,U)<\infty\)，
则称 \(E\) 在 \(\Omega\) 内具有局部有限周长。
后一条件下，\(D\mathbf1_E\) 是局部有限的向量 Radon 测度；对 Borel 集 \(A\subseteq\Omega\)，记
\[
 P(E,A)=|D\mathbf1_E|(A).
\]
全空间周长简记为 \(P(E)=P(E,\mathbb R^d)\)。
\end{definition}
\symidx{\(P(E,U)\)：集合在开集内的相对周长}
\symidx{\(P(E,\cdot)\)：有限周长集合的周长测度}

由变差对偶，测度记号在开集上与前面的定义一致。
示性函数总是局部可积，但在无界区域上未必属于 \(L^1(\Omega)\)。
因此，有限周长不自动意味着 \(\mathbf1_E\in BV(\Omega)\)；还需要 \(m_d(E)<\infty\)。
例如 \(E=\Omega\) 的相对周长为零，无论 \(\Omega\) 的体积是否有限。
相对周长也不计入位于 \(\partial\Omega\) 上的部分，不能与全空间周长混用。
补集与原集合具有相同的相对周长，因为
\(D\mathbf1_{\Omega\setminus E}=-D\mathbf1_E\) 于 \(\Omega\)。

\begin{lemma}[层集周长的可测性]\label{b1:lem:level-perimeter-measurable}
若 \(u\) 在 \(\Omega\) 上可测，则对每个开集 \(U\subseteq\Omega\)，
函数 \(t\mapsto P(E_t,U)\) 可测，允许取值 \(+\infty\)。
\end{lemma}
\begin{proof}
用紧集 \(K_j\) 穷竭 \(U\)，并要求每个紧支集测试场的支集最终包含在某个 \(K_j\) 内。
固定 \(j\)，考虑支集包含于 \(K_j\)、且 \(\|\varphi\|_\infty\le1\) 的
\(C_c^1\) 测试场。以函数及一阶导数的上确界之和作距离，这个集合可分：
将场及其全部一阶导数限制到 \(K_j\)，它嵌入有限个 \(C(K_j)\) 的乘积，
而紧度量空间上的连续函数空间可分。于是可以在这组测试场内部选一个可数稠密族。
将各 \(j\) 的稠密族合并，记为 \((\varphi_k)\)。

对支集固定的测试场，\(C^1\) 收敛使
\(\int_{E_t}\operatorname{div}\varphi\) 收敛，因为示性函数有界且支集体积有限。
因而对每个 \(t\)，都有
\[
 P(E_t,U)=\sup_k\int_U\mathbf1_{\{u(x)>t\}}
                         \operatorname{div}\varphi_k(x)\,\mathrm dx.
\]
每一项都是可测的含参变量积分：被积函数在共同的紧支集上绝对可积。
可数上确界仍可测，结论成立。
\end{proof}

\subsection{余面积公式}
\label{b1:subsec:270303}

\begin{lemma}[层集积分对变差的控制]\label{b1:lem:layer-perimeter-controls-variation}
设 \(u\in L^1_{\mathrm{loc}}(\Omega)\)。对任意开集 \(U\subseteq\Omega\)，有
\[
 V(u,U)\le\int_{\mathbb R}P(E_t,U)\,\mathrm dt.
\]
\end{lemma}
\begin{proof}
任取 \(\varphi\in C_c^1(U;\mathbb R^d)\)，\(\|\varphi\|_\infty\le1\)。
由\eqref{b1:eq:bv-signed-layer-cake}，且
\[
 \int_{\mathbb R}\int_U
 \bigl|\mathbf1_{\{u>t\}}-\mathbf1_{\{0>t\}}\bigr|
                 \cdot|\operatorname{div}\varphi|\,\mathrm dx\,\mathrm dt
 =\int_U|u|\cdot|\operatorname{div}\varphi|\,\mathrm dx<\infty,
\]
可以使用 Fubini 定理。紧支集又保证
\(\int_U\operatorname{div}\varphi\,\mathrm dx=0\)，故
\[
 \int_Uu\,\operatorname{div}\varphi\,\mathrm dx
 =\int_{\mathbb R}\left(\int_{E_t}\operatorname{div}\varphi\,\mathrm dx\right)\,\mathrm dt
 \le\int_{\mathbb R}P(E_t,U)\,\mathrm dt.
\]
先保留基准项验证绝对可积，再利用散度积分为零将它消去，是有符号分层的关键。
对测试场取上确界即得结论。
\end{proof}

\begin{lemma}[\(W^{1,1}\) 的层集周长公式]\label{b1:lem:w11-perimeter-coarea}
设 \(U\) 为有界开集，\(v\in W^{1,1}(U)\)。则
\[
 \int_{\mathbb R}P(\{v>t\},U)\,\mathrm dt
       =\int_U|\nabla v|\,\mathrm dx.
\]
\end{lemma}
\begin{proof}
取 \(\varepsilon>0\)，定义分段线性截断
\[
 H_\varepsilon(s)=\min\bigl\{1,\max\{0,s/\varepsilon\}\bigr\}.
\]
Sobolev 截断公式及水平集上的梯度为零给出
\[
 H_\varepsilon(v-t)\in W^{1,1}(U),\quad
 \bigl|\nabla H_\varepsilon(v-t)\bigr|
   =\varepsilon^{-1}\mathbf1_{\{t<v<t+\varepsilon\}}|\nabla v|
       \quad\text{几乎处处}.
\]
在固定的层值 \(t\) 上，\(H_\varepsilon(v-t)\) 逐点收敛到
\(\mathbf1_{\{v>t\}}\)，且值在 \([0,1]\) 内。
由于 \(U\) 有限体积，控制收敛给出 \(L^1(U)\) 收敛。
变差下半连续性于是给出
\[
 P(\{v>t\},U)
 \le\liminf_{\varepsilon\downarrow0}
      \varepsilon^{-1}\int_{\{t<v<t+\varepsilon\}}|\nabla v|\,\mathrm dx.
\]
沿任意趋零的正数列应用 Fatou 引理，再用 Tonelli 定理，得到
\[
 \begin{aligned}
 \int_{\mathbb R}P(\{v>t\},U)\,\mathrm dt
 &\le\liminf_{\varepsilon\downarrow0}
   \int_U|\nabla v(x)|\left(
    \varepsilon^{-1}\int_{\mathbb R}
       \mathbf1_{\{t<v(x)<t+\varepsilon\}}\,\mathrm dt\right)\,\mathrm dx\\
 &=\int_U|\nabla v|\,\mathrm dx.
 \end{aligned}
\]
括号中的积分等于 \(1\)，因为固定 \(x\) 后，所允许的 \(t\) 构成长度为
\(\varepsilon\) 的区间。反向不等式由前一引理和
\(V(v,U)=\int_U|\nabla v|\) 得到。
\end{proof}

\begin{theorem}[BV 余面积公式]\label{b1:thm:bv-coarea}
设 \(u\in BV(\Omega)\) 为实值函数，\(E_t=\{u>t\}\)。
则几乎每个 \(E_t\) 在 \(\Omega\) 内具有有限周长，且对每个开集 \(U\subseteq\Omega\)，
\begin{equation}\label{b1:eq:bv-coarea-open}
 |Du|(U)=\int_{\mathbb R}P(E_t,U)\,\mathrm dt.
\end{equation}
在周长无限的零测层值集合上，将周长测度统一记为零，则对任意 Borel 集
\(A\subseteq\Omega\)，都有
\begin{equation}\label{b1:eq:bv-coarea-measure}
 |Du|(A)=\int_{\mathbb R}|D\mathbf1_{E_t}|(A)\,\mathrm dt.
\end{equation}
\end{theorem}
\begin{proof}
先设 \(U\subseteq\Omega\) 有界。将 \(u\) 限制到 \(U\)，用%
\cref{b1:thm:bv-strict-approximation}选
\(u_j\in C^\infty(U)\cap W^{1,1}(U)\)，使
\[
 u_j\to u\quad\text{于 }L^1(U),\quad
 \int_U|\nabla u_j|\,\mathrm dx\longrightarrow|Du|(U).
\]
按\eqref{b1:eq:bv-level-set-l1-distance}抽一个共同子列，使几乎每个层值都有
\(\mathbf1_{\{u_j>t\}}\to\mathbf1_{E_t}\) 于 \(L^1(U)\)。
变差下半连续性和 Fatou 引理给出
\[
 \begin{aligned}
 \int_{\mathbb R}P(E_t,U)\,\mathrm dt
 &\le\liminf_j\int_{\mathbb R}P(\{u_j>t\},U)\,\mathrm dt\\
 &=\lim_j\int_U|\nabla u_j|\,\mathrm dx=|Du|(U).
 \end{aligned}
\]
最后的等号使用前一引理。由层集积分控制变差的引理，又有反向不等式。

对一般开集 \(U\)，用递增有界开集 \(U_j\) 穷竭它。
每个紧支集测试场最终支承于某个 \(U_j\)，故对所有 \(t\)，
\[
 P(E_t,U)=\lim_jP(E_t,U_j).
\]
对周长一侧用单调收敛定理，对 \(|Du|\) 一侧用测度从下连续性，
便得到\eqref{b1:eq:bv-coarea-open}。
取 \(U=\Omega\)，右端有限，所以周长无限的层值构成一个 Lebesgue 零测集 \(Z\)。

还须将开集等式推广到 Borel 集，而不为每个集合另选一组例外层值。
对 \(t\notin Z\)，令 \(\mu_t=|D\mathbf1_{E_t}|\)，对 \(t\in Z\) 令 \(\mu_t=0\)。
这些都是 \(\Omega\) 上的有限 Borel 测度，且对每个开集 \(U\)，
\(t\mapsto\mu_t(U)\) 可测。
使 \(t\mapsto\mu_t(A)\) 可测的 Borel 集 \(A\) 构成一个 Dynkin 系统：
补集利用 \(\mu_t(\Omega)<\infty\) 作差，可数不交并利用单调收敛。
它包含开集，而开集构成生成 Borel 集的 \(\pi\)-系统，故包含全部 Borel 集。
因此
\[
 \lambda(A)=\int_{\mathbb R}\mu_t(A)\,\mathrm dt
\]
定义一个有限 Borel 测度；可数可加性由 Tonelli 定理给出。
开集等式说明 \(\lambda\) 与 \(|Du|\) 在所有开集上一致，
有限测度唯一性遂给出它们在全部 Borel 集上一致。
这就是\eqref{b1:eq:bv-coarea-measure}。
\end{proof}

周长本身始终用分布导数定义。在光滑情形，它与通常边界面积的联系还需相应的
Gauss 公式；上述论证没有将这个联系预先放入定义。
对于一般 BV 函数，几乎所有层集具有有限周长并不意味着每个层集都有光滑边界，
也不意味着例外层值不存在。

\subsection{Cavalieri 型公式}
\label{b1:subsec:270304}

普通的 Cavalieri 分层公式分解体积，BV 余面积公式分解变差。
例如对 \(u\in BV(\Omega)\)，前者给出
\[
 \|u\|_{L^1(\Omega)}
 =\int_0^\infty m_d(\{u>t\})\,\mathrm dt
    +\int_0^\infty m_d(\{u<-t\})\,\mathrm dt,
\]
而变差等式对全部实数层值积分，不能在有符号情形只留下正层值。
例如取 \(u=c\mathbf1_E\)，其中 \(\mathbf1_E\in BV(\Omega)\)、\(c\ne0\)。
若 \(c>0\)，有周长贡献的层位于 \((0,c)\)，这些层集均为 \(E\)；
若 \(c<0\)，对应区间为 \((c,0)\)，层集则为 \(\Omega\setminus E\)。
补集周长不变，两个方向都给出 \(|Du|=|c|\cdot|D\mathbf1_E|\)。

由\eqref{b1:eq:bv-coarea-measure}，对每个非负 Borel 函数 \(g\)，还有加权形式
\begin{equation}\label{b1:eq:bv-weighted-coarea}
 \int_\Omega g\,\mathrm d|Du|
 =\int_{\mathbb R}\left(\int_\Omega g\,\mathrm d|D\mathbf1_{E_t}|\right)\,\mathrm dt.
\end{equation}
先对示性函数和非负简单函数应用测度等式，再作单调逼近，即可同时得到内层积分的可测性与该公式。
取 \(g=\mathbf1_A\)，它描述指定区域内的变差来自哪些层；
取一般权重，则可以分别强调不同空间位置。

分层还精确描述了截断所保留的变差。固定 \(a<b\)，令
\[
 F_{a,b}(s)=\min\bigl\{\max\{s-a,0\},b-a\bigr\}
                -\min\bigl\{\max\{-a,0\},b-a\bigr\}.
\]
减去常数使 \(F_{a,b}(0)=0\)，所以
\(|F_{a,b}(u)|\le|u|\)，在无限体积区域上也保留 \(L^1\) 可积性。

\begin{proposition}[截断的变差分层]\label{b1:prop:bv-truncation-coarea}
若 \(u\in BV(\Omega)\)，则 \(F_{a,b}(u)\in BV(\Omega)\)，且对每个 Borel 集
\(A\subseteq\Omega\)，
\[
 |D F_{a,b}(u)|(A)=\int_a^b|D\mathbf1_{E_t}|(A)\,\mathrm dt.
\]
\end{proposition}
\begin{proof}
截断有有符号表示
\[
 F_{a,b}(u)=\int_a^b
      \bigl(\mathbf1_{\{u>t\}}-\mathbf1_{\{0>t\}}\bigr)\,\mathrm dt.
\]
按层集控制变差引理的测试计算，
\[
 V(F_{a,b}(u),\Omega)\le\int_a^bP(E_t,\Omega)\,\mathrm dt
       \le|Du|(\Omega).
\]
结合 \(L^1\) 可积性，变差对偶的刻画给出 \(F_{a,b}(u)\in BV(\Omega)\)。
记 \(c=\min\bigl\{\max\{-a,0\},b-a\bigr\}\)。
当 \(-c<s<b-a-c\) 时，
\[
 \{F_{a,b}(u)>s\}=\{u>a+c+s\}.
\]
区间以外的超水平集为全区域或空集，只有两个端点层值不必遵循这段表述，
而它们不影响层积分。将这一关系代入 \(F_{a,b}(u)\) 的测度余面积公式，
再作平移换元 \(t=a+c+s\)，即得结论。
\end{proof}

因此，把函数值限制在一段区间内，恰好保留这段层值所贡献的变差。
这种按值域而非按空间切分的方法，将在有限周长集合的逼近以及后续精细结构中继续使用。

% !TeX root = ../../Book1.tex
\section{有限周长集合}
\label{b1:sec:2704}

前一节用示性函数的分布导数定义了周长。这一定义不需要先知道边界是否光滑，甚至不要求集合有有限体积。现在的问题是：导数测度集中在哪里，它的方向能否解释为边界法向？为此必须区分普通拓扑边界、按体积密度定义的边界，以及能从周长测度读出单位方向的约化边界。

\ProseParagraphBreak

本节固定开集 \(\Omega\subset\mathbb R^d\)，令 \(E\subset\mathbb R^d\) 为在 \(\Omega\) 内局部有限周长的可测集。所有边界与支撑均只考察 \(\Omega\) 内的点。需要时可先将 \(E\) 换成几乎处处相等的 Borel 集；下面的对象不会因此改变。

% 实读本地 Simon1983Lectures，正式第14节印刷71--76页；
% 本节特别对应72页(14.1)--(14.2)及极分解后的微分论证。
% 原书用内法向表示Dχ_E/|Dχ_E|；本书外法向取其负号。
% 本节不提前把约化边界认作Hausdorff面积边界，几何识别在下一节证明。

\subsection{周长测度}
\label{b1:subsec:270401}

沿用\cref{b1:def:perimeter}，记
\[
 \mu_E=|D\chi_E|,\quad P(E,A)=\mu_E(A)
\]
对 Borel 集 \(A\subset\Omega\) 成立。局部有限性是说每个 \(K\Subset\Omega\) 上的 \(\mu_E(K)\) 有限，不等于 \(P(E,\Omega)<\infty\)，更不等于 \(m_d(E)<\infty\)。例如半空间的体积与全空间周长都无穷，但在每个有界内部区域中的周长有限。

\ProseParagraphBreak

周长只依赖集合的几乎处处类。若 \(m_d(E\mathbin{\triangle}F)=0\)，则 \(\chi_E=\chi_F\) 几乎处处，所以它们的分布导数及全变差完全相同。对于补集，内部测试函数的积分给出
\begin{equation}\label{b1:eq:perimeter-complement}
 D\chi_{\mathbb R^d\setminus E}=-D\chi_E,\quad
 \mu_{\mathbb R^d\setminus E}=\mu_E.
\end{equation}
这说明周长本身不区分内外，区分内外的是导数测度的方向。

\begin{proposition}\label{b1:prop:perimeter-support-locality}
对 \(x\in\Omega\)，下列条件等价：\(x\in\operatorname{supp}\mu_E\)；每个满足 \(\overline B(x,r)\Subset\Omega\) 的球都有
\[
 0<m_d\bigl(E\cap B(x,r)\bigr)<m_d\bigl(B(x,r)\bigr).
\]
\end{proposition}
\begin{proof}
若某个球内 \(\chi_E\) 几乎处处为 \(0\) 或 \(1\)，则内部所有分布偏导为零，周长测度在该球内也为零，故 \(x\) 不在支撑中。

反过来，若 \(\mu_E\bigl(B(x,r)\bigr)=0\)，则 \(\chi_E\) 在此球中的所有分布偏导为零。局部磨光后，所得光滑函数在每个较小的同心球上梯度为零，因而为常数。让磨光尺度趋零，常数的 \(L^1\) 极限说明 \(\chi_E\) 在每个较小球上几乎处处为常数；重叠处这些常数相同。由于示性函数只取 \(0,1\)，它在原球内也几乎处处取其中一个值。
\end{proof}

这个判据只说明每个球都看得到两侧的正体积，还没有给出两侧所占的比例。某一侧的比例可能随半径缩小而趋于零；不能把周长的支撑直接定义成下面的测度论边界。

\subsection{测度论内部与外部}
\label{b1:subsec:270402}

\begin{definition}\label{b1:def:measure-interior-exterior}
给定可测集 \(E\)，定义它在 \(\Omega\) 内的\term[cedulun-neibu]{测度论内部}[measure-theoretic interior]与\term[cedulun-waibu]{测度论外部}[measure-theoretic exterior]分别为
\[
 E^1\coloneqq\left\{x\in\Omega:
 \lim_{r\downarrow0}\frac{m_d\bigl(E\cap B(x,r)\bigr)}{\omega_dr^d}=1\right\},
\]
\[
 E^0\coloneqq\left\{x\in\Omega:
 \lim_{r\downarrow0}\frac{m_d\bigl(E\cap B(x,r)\bigr)}{\omega_dr^d}=0\right\}.
\]
\end{definition}
\symidx{\(E^1\)：集合的测度论内部}
\symidx{\(E^0\)：集合的测度论外部}

这是\chapref{b1:ch:14}密度点概念的集合记号，而不是另一种拓扑内部。球体积积分的 Borel 性及有理半径逼近说明 \(E^1,E^0\) 都是 Borel 集。\cref{b1:thm:lebesgue-density}给出
\[
 m_d\bigl((E\cap\Omega)\mathbin{\triangle}E^1\bigr)=0,\quad
 m_d\bigl((\Omega\setminus E)\mathbin{\triangle}E^0\bigr)=0.
\]
因此 \(E^1\) 是集合的一个由密度选定的代表，但它未必开。普通内点属于 \(E^1\)，普通外点属于 \(E^0\)；逆命题没有普遍保证。

\ProseParagraphBreak

例如在 \(\Omega=\mathbb R^d\) 中取 \(E=B(0,1)\setminus\mathbb Q^d\)。这个集合没有普通内点，却有 \(E^1=B(0,1)\)；它与开球具有完全相同的示性函数等价类和周长。若使用任意代表的拓扑边界去计算周长，就会把整个闭球误当成边界。

\subsection{测度论边界}
\label{b1:subsec:270403}

\begin{definition}\label{b1:def:measure-theoretic-boundary}
给定可测集 \(E\)，称
\[
 \partial^eE\coloneqq\Omega\setminus(E^0\cup E^1)
\]
为它在 \(\Omega\) 内的\term[cedulun-bianjie]{测度论边界}[measure-theoretic boundary]，也在外文索引中以 essential boundary 登记。\foreignidx{essential boundary}
\end{definition}
\symidx{\(\partial^eE\)：测度论边界}

定义允许密度不存在。因此 \(\partial^eE\) 并不只是密度介于 \(0\) 与 \(1\) 之间且极限存在的点集。它是 Borel 集，并且对任意可测 \(E\) 都有 \(m_d(\partial^eE)=0\)。但在余维一的尺度上，这个集合可以很大。

\ProseParagraphBreak

半空间边界上每个球被平面分成等体积的两半，故其边界点的密度为 \(1/2\)。对于第一象限 \(E=(0,\infty)^2\)，正坐标轴上的非原点边界点仍有密度 \(1/2\)，原点的密度却为 \(1/4\)。它们都属于测度论边界；仅知道某点在 \(\partial^eE\)，还不能给它指定唯一切线或法向。

\ProseParagraphBreak

对于任意选定的集合代表，\(\partial^eE\) 包含于它的普通边界在 \(\Omega\) 中的部分：普通内外点都已经被密度 \(1\) 或 \(0\) 排除。前述删去有理点的例子表明，这个包含可以非常严格。与普通边界不同，\(\partial^eE\) 在零测修改下保持不变。

\subsection{约化边界}
\label{b1:subsec:270404}

由\cref{b1:prop:vector-measure-polar}，写
\[
 D\chi_E=\sigma_E\mu_E,\quad |\sigma_E|=1
 \quad\mu_E\text{-几乎处处}.
\]
一个点附近的方向可能彼此抵消。只有当缩小球中方向的平均趋于单位向量时，才看得到一个没有被抵消的极限方向。

\begin{definition}\label{b1:def:reduced-boundary}
设 \(E\) 在 \(\Omega\) 内局部有限周长。若 \(x\in\operatorname{supp}\mu_E\)，并且极限
\begin{equation}\label{b1:eq:reduced-boundary-direction}
 a_E(x)\coloneqq\lim_{r\downarrow0}
 \frac{D\chi_E\bigl(\overline B(x,r)\bigr)}
 {\mu_E\bigl(\overline B(x,r)\bigr)}
\end{equation}
存在且长度为 \(1\)，则称 \(x\) 为 \(E\) 的约化边界点。全体这些点组成\term[yuehua-bianjie]{约化边界}[reduced boundary]，记为 \(\partial^*E\)。
\end{definition}
\symidx{\(\partial^*E\)：由周长方向极限定义的约化边界}

这里 \(\partial^*E\) 专指约化边界，\(\partial^eE\) 专指测度论边界；不能因文献记号不同而互换二者。式\eqref{b1:eq:reduced-boundary-direction}采用与 Radon 微分定理相同的闭球。将半径从内增大或从外减小，并对向量测度的各分量使用连续性，可知用开球得到的极限完全相同。在支撑上，两种球的分母在充分小的正半径下都为正。因此，这一定义与几何测度论文献中通常采用开球的约化边界定义等价，并未引入另一种边界概念。

\begin{proposition}\label{b1:prop:reduced-boundary-carries-perimeter}
约化边界为 Borel 集，\(a_E\) 可取为其上的 Borel 函数，并且
\[
 \mu_E(\Omega\setminus\partial^*E)=0.
\]
对每个 \(x\in\partial^*E\)，还有
\begin{equation}\label{b1:eq:reduced-direction-oscillation}
 \lim_{r\downarrow0}
 \frac{1}{\mu_E\bigl(B(x,r)\bigr)}
 \int_{B(x,r)}|\sigma_E(y)-a_E(x)|^2\,\mathrm d\mu_E(y)=0.
\end{equation}
\end{proposition}
\begin{proof}
球质量与向量分量的球积分都是 Borel 函数。对固定中心，将闭球半径从外作有理数逼近，利用局部有限测度的连续性，可把极限存在及极限长度为 \(1\) 的条件写成可数个 Borel 条件；支撑本身闭，故约化边界 Borel，极限函数也 Borel。

由\cref{b1:thm:radon-function-differentiation}逐分量作用于 \(\sigma_E\)，球平均在 \(\mu_E\)-几乎处处趋于 \(\sigma_E(x)\)。后者长度为 \(1\)，而支撑的补集为 \(\mu_E\)-零集，故周长集中在约化边界上。最后，固定该边界的一点，利用两个方向的单位长度得到
\[
 \frac{\int_{B(x,r)}|\sigma_E-a_E(x)|^2\,\mathrm d\mu_E}
 {\mu_E\bigl(B(x,r)\bigr)}
 =2-2a_E(x)\cdot
 \frac{D\chi_E\bigl(B(x,r)\bigr)}{\mu_E\bigl(B(x,r)\bigr)},
\]
右端趋于零，证明了最后一式。
\end{proof}

这一命题仍是向量测度的微分结论，没有证明 \(\mu_E\) 等于某种 Hausdorff 面积。下一节的结构定理将说明：约化边界处的集合缩放为半空间，周长缩放为其边界平面的面积测度。

\subsection{测度论法向}
\label{b1:subsec:270405}

\begin{definition}\label{b1:def:measure-theoretic-outer-normal}
对 \(x\in\partial^*E\)，定义 \(E\) 的\term[cedulun-waifaxiang]{测度论外法向}[measure-theoretic outer unit normal]为
\[
 \nu_E(x)\coloneqq-a_E(x).
\]
\end{definition}
\symidx{\(\nu_E(x)\)：有限周长集合的测度论外法向}

负号来自分布导数的约定。例如 \(E=\{x_d<0\}\) 时，直接沿 \(x_d\) 积分得到
\[
 \int_E\partial_d\varphi\,\mathrm dx
 =\int_{\mathbb R^{d-1}}\varphi(x',0)\,\mathrm dx',
\]
所以 \(D\chi_E=-e_d\mathcal H^{d-1}\llcorner\{x_d=0\}\)，而几何上的外法向是 \(e_d\)。一般情形因此写成
\[
 D\chi_E=-\nu_E\mu_E.
\]
法向只需在 \(\partial^*E\) 上指定；在其补集任意补值不会改变测度等式。补集的约化边界与原集相同，而外法向相反，这也与\eqref{b1:eq:perimeter-complement}一致。

\ProseParagraphBreak

第一象限的角点说明单位长度条件有实际作用。由两次一维分部积分，
\[
 D\chi_{(0,\infty)^2}
 =e_1\mathcal H^1\llcorner\bigl(\{0\}\times(0,\infty)\bigr)
 +e_2\mathcal H^1\llcorner\bigl((0,\infty)\times\{0\}\bigr).
\]
原点半径 \(r\) 的球中，两条射线各贡献长度 \(r\)，所以向量与全变差之比为 \((e_1+e_2)/2\)，长度为 \(1/\sqrt2\)。原点属于 \(\partial^eE\)，却不属于 \(\partial^*E\)。这不是在角点随意漏掉一个法向，而是定义准确地记录了两个方向的抵消。

\ProseParagraphBreak

Simon 的《Lectures on Geometric Measure Theory》以 \(D\chi_E/\mu_E\) 表示内法向\cite[\S14, 14.1--14.2]{Simon1983Lectures}。本书选择外法向，是为使下一节的 Gauss–Green 公式保持“区域内散度积分等于向外通量”的符号。阅读其他文献时，应先检查这一等式，再对照法向名称。

% !TeX root = ../../Book1.tex
\section{Gauss–Green 与 De Giorgi 结构}
\label{b1:sec:2705}

上一节已经得到带方向的周长测度，但尚未证明它是一份边界面积。本节把两件事接起来。先对光滑图域直接计算向外通量，再在约化边界点缩放集合；紧性给出极限，方向的集中使极限只能是半空间，而两侧体积下界使其边界必须穿过缩放中心。最后从这些半空间逼近中构造可数张 Lipschitz 图，并辨认周长的密度。

% 实读 Simon1983Lectures，正式第14节印刷71--76页：
% 72--73页(14.1)--(14.3)及定理；73--76页约化点方向、周长上界、
% BV紧性、半空间极限与两侧体积下界的完整证明。
% 本书外法向=-Dχ/|Dχ|，统一把法向e_d对应的极限写成{x_d<0}。
% 原书引用11.6与3.5完成可整流/测度识别；本稿改为显式锥排除图分片，
% 再调用本书14章Radon微分及19章已完整证明的带权密度命题。
% 相对等周用22章一阶球延拓、24章GNS与27.2严格逼近自足推出；
% ∂^eE\∂*E的Hausdorff零性另外给5r证明，不冒称Simon14单独陈述了全部扩充结论。

\subsection{Gauss–Green 公式}
\label{b1:subsec:270501}

先把分布恒等式写成\term*[gauss-green-gongshi]{Gauss--Green 公式}[Gauss--Green formula]。它此时涉及周长测度，面积形式还需要后面的几何识别。

\begin{proposition}[分布形式的 Gauss--Green 公式]\label{b1:prop:gauss-green-distributional}
设 \(E\) 在 \(\Omega\) 中局部有限周长，\(\Phi\in C_c^1(\Omega;\mathbb R^d)\)。则
\begin{equation}\label{b1:eq:gauss-green-perimeter-measure}
 \int_E\operatorname{div}\Phi\,\mathrm dx
 =\int_{\partial^*E}\Phi\cdot\nu_E\,\mathrm d\mu_E.
\end{equation}
\end{proposition}
\begin{proof}
分布导数的定义给出左侧等于 \(-\int\Phi\cdot\mathrm dD\chi_E\)。代入 \(D\chi_E=-\nu_E\mu_E\)，再用\cref{b1:prop:reduced-boundary-carries-perimeter}，即得公式。最初的光滑测试公式可以通过固定紧支集上的 \(C^1\) 逼近传给 \(C_c^1\) 向量场。
\end{proof}

此时边界积分仍是对 \(\mu_E\) 积分，不是已经知道的 Hausdorff 面积积分。对光滑边界，可以不依赖结构定理而直接识别它。

\begin{proposition}[光滑区域的周长]\label{b1:prop:smooth-domain-perimeter}
设开集 \(E\) 在 \(\Omega\) 内具有单侧 \(C^1\) 图边界：每个边界点附近，适当刚性变换后 \(E\) 都是某个 \(C^1\) 函数图像的一侧。则 \(E\) 在 \(\Omega\) 中局部有限周长，而且
\begin{equation}\label{b1:eq:smooth-domain-perimeter}
 D\chi_E=-\nu_E\mathcal H^{d-1}\llcorner(\partial E\cap\Omega),
 \quad
 P(E,A)=\mathcal H^{d-1}(\partial E\cap A)
\end{equation}
对 Borel 集 \(A\subset\Omega\) 成立。这里 \(\nu_E\) 是通常的外单位法向，且 \(\partial^*E=\partial E\cap\Omega\)。
\end{proposition}
\begin{proof}
先设 \(d\ge2\)，并在一个图表内写 \(x=(z,t)\)、\(E=\{t<g(z)\}\)，其中 \(g\in C^1\)。向量场支集紧含图表，所以沿竖直方向积分时，图表上下端没有额外项。微积分基本定理给出
\[
 \int_E\partial_d\Phi_d\,\mathrm dx
 =\int\Phi_d\bigl(z,g(z)\bigr)\,\mathrm dz.
\]
对 \(i<d\)，先对
\(\int_{-\infty}^{g(z)}\Phi_i(z,t)\,\mathrm dt\)
作 \(z_i\) 微分，再对 \(z\) 积分，得到
\[
 \int_E\partial_i\Phi_i\,\mathrm dx
 =-\int\Phi_i\bigl(z,g(z)\bigr)\,\partial_i g(z)\,\mathrm dz.
\]
因而图表内的散度积分等于
\[
 \int \Phi\bigl(z,g(z)\bigr)\cdot(-\nabla g(z),1)\,\mathrm dz.
\]
\cref{b1:thm:graph-area-formula}说明图面积元为
\(\sqrt{1+|\nabla g|^2}\,\mathrm dz\)，而
\[
 \nu_E\bigl(z,g(z)\bigr)
 =\frac{(-\nabla g(z),1)}{\sqrt{1+|\nabla g(z)|^2}}.
\]
这就得到局部公式。

对于一般紧支集测试场，以有限个边界图表及不交边界的内部区域覆盖其支集，再取光滑单位分解。分别计算各个分片后相加，分解系数的梯度项因其和恒为 \(1\) 而抵消。在不交边界的区域，示性函数局部常数，散度积分为零。每个紧集附近只需有限张图，其面积有限，所以所得公式同时证明局部有限周长。

向量测度表示的唯一性给出 \(D\chi_E\) 的等式；法向长度为 \(1\)，故其全变差就是边界面积。法向在每张图上连续，缩小球中的平均趋于该点的法向，于是每个光滑边界点都约化；内外点则不在周长支撑中。一维情形直接对区间使用微积分基本定理，每个内部端点贡献一个单位点质量，结论相同。
\end{proof}

“单侧”不能删除。如果只要求拓扑边界恰为一张光滑曲面，却允许集合占据曲面的两侧，例如删去一张超平面后的全空间，那么示性函数几乎处处恒为 \(1\)，周长实际上为零。

特别地，对球 \(E=B(a,R)\)，有
\[
 \nu_E(x)=\frac{x-a}{R}\quad(x\in\partial B(a,R)),\quad
 P\bigl(B(a,R),A\bigr)=\mathcal H^{d-1}\bigl(\partial B(a,R)\cap A\bigr).
\]
这个计算已经可以独立使用，不必等待一般有限周长集合的结构。

\subsection{约化边界的可整流性}
\label{b1:subsec:270502}

下面的证明先在 \(d\ge2\) 中进行。一维的结构将在定理后单独说明。第一个工具是球内的\term[xiangdui-dengzhou-budengshi]{相对等周不等式}[relative isoperimetric inequality]：如果球内两侧都有一定体积，它们之间就必须有一定周长。

\begin{lemma}[球内相对等周不等式]\label{b1:lem:relative-isoperimetric-ball}
设 \(d\ge2\)，\(\overline B(x,r)\Subset\Omega\)，且 \(E\) 在 \(\Omega\) 内局部有限周长。则
\begin{equation}\label{b1:eq:relative-isoperimetric-ball}
 \min\Bigl(m_d\bigl(E\cap B(x,r)\bigr),m_d\bigl(B(x,r)\setminus E\bigr)\Bigr)^{(d-1)/d}
 \le C_d\mu_E\bigl(B(x,r)\bigr).
\end{equation}
\end{lemma}
\begin{proof}
先对单位球内的 \(v\in W^{1,1}\) 证明
\[
 \|v-v_B\|_{L^{d/(d-1)}(B)}
 \le C_d\|\nabla v\|_{L^1(B)}.
\]
由\cref{b1:thm:ball-poincare-lp}，\(v-v_B\) 的 \(L^1\) 范数受梯度控制。用\cref{b1:thm:lipschitz-domain-extension}把它延拓到全空间，其梯度 \(L^1\) 范数仍有同一控制；再用\cref{b1:thm:sobolev-subcritical-embedding}的 \(p=1\) 情形即得结论。伸缩到任意球时，左侧与梯度积分都按 \(r^{d-1}\) 变化，故常数不依赖球。

将\cref{b1:thm:bv-strict-approximation}用于球中的 \(\chi_E\)。逼近函数在 \(L^1\) 中收敛，其均值也收敛，而梯度积分趋于 \(\mu_E(B)\)。抽取几乎处处收敛子列，用 Fatou 引理得到
\[
 \|\chi_E-a\|_{L^{d/(d-1)}(B)}
 \le C_d\mu_E(B),\quad a=\frac{m_d(E\cap B)}{m_d(B)}.
\]
若 \(a\le1/2\)，左侧至少为
\(\tfrac12m_d(E\cap B)^{(d-1)/d}\)；若 \(a\ge1/2\)，则改在补集上估计。这就证明了相对等周不等式。
\end{proof}

\begin{theorem}[约化边界处的半空间缩放]\label{b1:thm:reduced-boundary-blowup}
设 \(d\ge2\)，\(x\in\partial^*E\)，并记
\[
 E_{x,r}=\{z:x+rz\in E\},\quad
 H_x=\{z:z\cdot\nu_E(x)<0\}.
\]
则当 \(r\downarrow0\) 时，
\begin{equation}\label{b1:eq:reduced-boundary-halfspace}
 \chi_{E_{x,r}}\longrightarrow\chi_{H_x}
 \quad\text{于 }L^1_{\mathrm{loc}}(\mathbb R^d).
\end{equation}
而且对每个 \(\varphi\in C_c(\mathbb R^d)\)，
\begin{equation}\label{b1:eq:perimeter-flat-blowup}
 r^{1-d}\int_\Omega
 \varphi\left(\frac{y-x}{r}\right)\,\mathrm d\mu_E(y)
 \longrightarrow
 \int_{\nu_E(x)^\perp}\varphi\,\mathrm d\mathcal H^{d-1}.
\end{equation}
特别地，每个约化边界点都有集合密度 \(1/2\)，并且
\begin{equation}\label{b1:eq:reduced-perimeter-unit-density}
 \lim_{r\downarrow0}
 \frac{\mu_E\bigl(B(x,r)\bigr)}{\omega_{d-1}r^{d-1}}=1.
\end{equation}
\end{theorem}
\begin{proof}
经平移、旋转，设 \(x=0\)、\(\nu_E(0)=e_d\)。因此极方向的极限为 \(-e_d\)。以下简记 \(\mu=\mu_E\)、\(\sigma=\sigma_E\)，以及
\[
 V(r)=m_d\bigl(E\cap B(0,r)\bigr),\quad
 W(r)=m_d\bigl(B(0,r)\setminus E\bigr).
\]
球壳的体积界说明 \(V,W\) 绝对连续，且它们的导数几乎处处非负、不超过 \(d\omega_dr^{d-1}\)。方向极限给出：对所有充分小的半径，
\[
 \int_{B(0,r)}(-\sigma_d)\,\mathrm d\mu
 \ge\tfrac12\mu\bigl(B(0,r)\bigr).
\]

取在 \(B(0,r)\) 上为 \(1\)、在球壳 \(r<|y|<r+h\) 中线性降到零的径向截断，并用光滑截断逼近它。分布公式给出
\[
 \int(-\sigma_d)\zeta\,\mathrm d\mu
 =\int_E\partial_d\zeta\,\mathrm dx
 \le h^{-1}\bigl(V(r+h)-V(r)\bigr).
\]
在导数存在且球面不带 \(\mu\) 质量的半径令 \(h\downarrow0\)。这些半径占满一维测度，故得到
\begin{equation}\label{b1:eq:reduced-perimeter-radial-bound}
 \mu\bigl(B(0,r)\bigr)\le2V'(r)\le2d\omega_dr^{d-1}
 \quad\text{对几乎每个充分小的 }r.
\end{equation}
补集的外法向相反，同样有 \(\mu\bigl(B(0,r)\bigr)\le2W'(r)\)。由球的单调性，从较大半径向下逼近，还得到
\(\mu\bigl(B(0,r)\bigr)\le C_dr^{d-1}\)
对每个充分小的 \(r\) 成立。

现在排除退化的缩放极限。令 \(M(r)=\min(V(r),W(r))\)。由\cref{b1:prop:perimeter-support-locality}，\(M(r)>0\)；它绝对连续且 \(M(r)\to0\)。结合相对等周估计与刚才的两个导数界，得到
\[
 M(r)^{(d-1)/d}
 \le C_d\mu\bigl(B(0,r)\bigr)
 \le C'_d\min(V'(r),W'(r))
 \le C'_d M'(r)
\]
几乎处处成立。因此 \((M^{1/d})'\ge c_d>0\) 几乎处处。从正半径积分后再令下端趋零，得到
\begin{equation}\label{b1:eq:reduced-two-sided-volume}
 V(r)\ge c_dr^d,\quad W(r)\ge c_dr^d
\end{equation}
对全部充分小的 \(r\) 成立。这一步同时保证两侧不会消失。

对任意趋零序列 \(r_j\)，作变量代换可得
\[
 |D\chi_{E_{0,r_j}}|(B(0,R))
 =r_j^{1-d}\mu\bigl(B(0,r_jR)\bigr)\le C_dR^{d-1}
\]
对每个固定 \(R\) 及充分大的 \(j\) 成立。零阶 \(L^1\) 范数也受球体积控制。逐球使用\cref{b1:thm:bv-local-compactness}并对角抽列，得到
\(\chi_{E_{0,r_j}}\to\chi_F\)
于 \(L^1_{\mathrm{loc}}\)，其中 \(F\) 局部有限周长。这里各缩放定义域最终包含任意固定球；极限仍为示性函数，可由进一步的几乎处处收敛看出。

式\eqref{b1:eq:reduced-direction-oscillation}和 Cauchy 不等式说明
\[
 r^{1-d}\int_{B(0,rR)}|\sigma+e_d|\,\mathrm d\mu
 \longrightarrow0
\]
对固定 \(R\) 成立。于是缩放后的分布导数在前 \(d-1\) 个方向趋于零，在第 \(d\) 个方向的极限是非正分布。故
\[
 \partial_i\chi_F=0\quad(i<d),\quad
 \partial_d\chi_F\le0.
\]
将 \(\chi_F\) 磨光：前一组等式使光滑函数只依赖最后一个变量，后一不等式使它对这个变量非增。让磨光尺度趋零，就得到 \(\chi_F(z,t)=h(t)\) 几乎处处，其中 \(h\) 有非增代表且几乎处处只取 \(0,1\)。所以 \(F\) 几乎处处等于 \(\{t<b\}\)，其中暂允许 \(b=\pm\infty\)。

对任意固定 \(s>0\)，将\eqref{b1:eq:reduced-two-sided-volume}伸缩并传给极限，球 \(B(0,s)\) 中 \(F\) 及其补集的体积都至少为 \(c_ds^d\)。若 \(b\ne0\)，取足够小的 \(s\)，其中一侧就会为空，矛盾。因此 \(b=0\)。任意初始半径序列都有子列趋于同一个半空间，这也证明了全部半径下的\eqref{b1:eq:reduced-boundary-halfspace}。

最后记缩放周长测度为
\(\mu_r(A)=r^{1-d}\mu(rA)\)。
对 \(\varphi\in C_c^1\)，方向误差趋零给出
\[
 \int\varphi\,\mathrm d\mu_r
 =-\int\varphi\,\mathrm d(D_d\chi_{E_{0,r}})+o(1)
 =\int_{E_{0,r}}\partial_d\varphi\,\mathrm dx+o(1)
 \longrightarrow\int_{\{t=0\}}\varphi\,\mathrm d\mathcal H^{d-1}.
\]
紧集上的一致质量界允许用光滑函数一致逼近一般 \(C_c\) 测试函数，得到\eqref{b1:eq:perimeter-flat-blowup}。平面测度不给单位球面质量，故用内外连续截断逼近球的示性函数，可得\eqref{b1:eq:reduced-perimeter-unit-density}。集合密度 \(1/2\) 则直接来自半空间的体积及 \(L^1\) 收敛。
\end{proof}

半空间极限已经给出法向的几何意义，但可数图分片仍须构造。下面把每点成立的极限整理为统一尺度上的估计。

\begin{proposition}\label{b1:prop:reduced-boundary-rectifiable}
当 \(d\ge2\) 时，\(\partial^*E\) 可由可数张完整 Lipschitz 图覆盖，因而是可数 \(d-1\) 维可整流 Borel 集。
\end{proposition}
\begin{proof}
令 \(c_d\) 为\eqref{b1:eq:reduced-two-sided-volume}中的统一常数，固定
\(0<\varepsilon<c_d/8^d\)。
按极限开始满足估计的尺度，把约化边界分成可数个 Borel 片，同时要求 \(B(x,1/m)\subset\Omega\)，使每片中各点 \(x\) 对 \(0<r<1/m\) 都满足两侧体积下界，以及
\[
 m_d\bigl((E\mathbin{\triangle}\{y:(y-x)\cdot\nu_E(x)<0\})
                   \cap B(x,r)\bigr)\le\varepsilon r^d.
\]
球体积函数对半径连续，故只检查有理半径便可保证这些片 Borel。再按法向所落入的球面小邻域分片，使一片上
\(|\nu_E(x)-\nu_0|<1/8\)，其中 \(\nu_0\) 固定；最后与直径小于 \(1/(4m)\) 的可数小立方体相交。

取同一片中不同的 \(x,y\)，记 \(\ell=|x-y|\)。若
\(|(y-x)\cdot\nu_E(x)|>\ell/2\)，则 \(B(y,\ell/4)\) 完全位于以 \(x\) 为中心的极限半空间的一侧，并且包含于 \(B(x,2\ell)\)。在 \(y\) 处的两侧体积下界说明，这个小球对半空间的体积误差至少为 \(c_d(\ell/4)^d\)；在 \(x\) 处的误差估计却使它至多为 \(\varepsilon(2\ell)^d\)，与 \(\varepsilon\) 的选择矛盾。因此
\[
 |(y-x)\cdot\nu_0|
 \le |(y-x)\cdot\nu_E(x)|+\tfrac18|y-x|
 \le\tfrac58|y-x|.
\]
投影到 \(\nu_0^\perp\) 后，距离至少为 \(\sqrt{39}\,|y-x|/8\)。所以投影在该片上单射，法向坐标是投影坐标的 Lipschitz 函数，常数至多为 \(5/\sqrt{39}\)。由\cref{b1:thm:mcshane-extension}将这个实值函数延拓到整个 \(\nu_0^\perp\)，便得到包含该片的一张完整 Lipschitz 图。所有分片均可数，证明完成。
\end{proof}

这个论证没有从“极限是平面”直接跳到可整流性。关键是用两侧体积下界排除同一片上沿法向分离过大的两个点，进而让固定的正交投影成为单射。它也没有调用\chapref{b1:ch:19}只给出证明概要的一般密度逆判据。

\subsection{De Giorgi 结构定理}
\label{b1:subsec:270503}

\term*[de-giorgi-jiegou-dingli]{De Giorgi 结构定理}[De Giorgi structure theorem]把上述几何逼近与周长测度的表示合在一起。

\begin{theorem}[De Giorgi]\label{b1:thm:de-giorgi-structure}
设 \(E\) 在开集 \(\Omega\subset\mathbb R^d\) 中局部有限周长。当 \(d\ge2\) 时，\(\partial^*E\) 可数 \(d-1\) 维可整流；当 \(d=1\) 时，它是局部有限的点集。两种情形均有
\begin{equation}\label{b1:eq:de-giorgi-perimeter-measure}
 |D\chi_E|=\mathcal H^{d-1}\llcorner\partial^*E.
\end{equation}
此外，
\begin{equation}\label{b1:eq:essential-reduced-boundary}
 \partial^*E\subset\partial^eE,\quad
 \mathcal H^{d-1}(\partial^eE\setminus\partial^*E)=0.
\end{equation}
每个 \(x\in\partial^*E\) 的集合密度为 \(1/2\)。当 \(d\ge2\) 时，在\chapref{b1:ch:19}的缩放测试意义下，
\[
 \operatorname{Tan}^{d-1}(\partial^*E,x)=\nu_E(x)^\perp.
\]
\end{theorem}
\begin{proof}
先设 \(d\ge2\)，记 \(S=\partial^*E\)、\(k=d-1\)、\(\mu=\mu_E\)。可整流性和单位周长密度已经证明，且 \(\mu\) 集中在 \(S\) 上。先证明 \(\mu\ll\mathcal H^k\llcorner S\)。把 \(S\) 分成可数个片，在每片上都有
\[
 \mu\bigl(B(x,r)\bigr)\le C_dr^k\quad(0<r<1/m),
\]
并限制在距 \(\Omega\) 外部有正距离的有界区域内。若片中的集合 \(N\) 为 \(\mathcal H^k\)-零集，用直径小于 \(1/(2m)\) 的集合覆盖它，使直径的 \(k\) 次和任意小。每个覆盖集若与 \(N\) 相交，就以交点为中心取半径为其直径两倍的球；这些球仍覆盖 \(N\)，上式控制其总 \(\mu\) 质量。直径为零的单点没有质量，也由该上界看出。于是 \(\mu(N)=0\)。对各片取并，即得绝对连续性。

可数图覆盖保证 \(\mathcal H^k\llcorner S\) 是 \(\sigma\) 有限测度，故 Radon–Nikodym 定理给出
\[
 \mu=\theta\mathcal H^k\llcorner S.
\]
由已完整证明的\cref{b1:prop:rectifiable-weight-density}，\(\mu\) 的 \(k\) 维密度在 \(\mathcal H^k\)-几乎处处等于 \(\theta\)。而\eqref{b1:eq:reduced-perimeter-unit-density}说明该密度在每个 \(S\) 中的点都为 \(1\)，所以 \(\theta=1\) 几乎处处，证明了测度表示。特别地，未加权的约化边界面积现在也已知局部有限。将测度表示代入\eqref{b1:eq:perimeter-flat-blowup}，正好得到\chapref{b1:ch:19}定义中的单位重数切平面，故其方向为 \(\nu_E(x)^\perp\)。

半空间极限已经证明 \(S\subset\partial^eE\)。为了处理余项，注意：若 \(x\in\partial^eE\)，则存在 \(\delta_x>0\) 及任意小的半径 \(r\)，使球内两侧的体积分数均不小于 \(\delta_x\)。否则，连续函数
\[
 r\longmapsto\frac{m_d\bigl(E\cap B(x,r)\bigr)}{\omega_dr^d}
\]
在足够小半径下总落在分别靠近 \(0\)、\(1\) 的两个不交区间内；连续性使它只能停留在其中一侧，最终便趋于 \(0\) 或 \(1\)，与 \(x\in\partial^eE\) 矛盾。相对等周不等式于是给出任意小半径上的
\[
 \mu\bigl(B(x,r)\bigr)\ge c_xr^k,\quad c_x>0.
\]
将 \(N=\partial^eE\setminus S\) 按
\(\limsup_{r\downarrow0}\mu\bigl(B(x,r)\bigr)/r^k>1/j\)
分片，并与内部紧穷竭相交。球质量的 Borel 性及有理半径逼近说明所得各片 Borel；每片上都有任意小半径满足 \(\mu\bigl(B(x,r)\bigr)\ge r^k/j\)。固定其中一片 \(N_j\)，它是 \(\mu\)-零集。给定 \(\varepsilon>0\)，在一个内部有限质量的邻域中选取包含它的开集 \(G\)，使 \(\mu(G)<\varepsilon\)，再对每个点选取任意小、包含于 \(G\) 且满足上述下界的球。由\cref{b1:lem:five-r-covering}选出不交球 \(B_i=B(x_i,r_i)\)，其五倍球覆盖该片。因此在任意小的 Hausdorff 覆盖尺度上，
\[
 \mathcal H^k_{\delta}(N_j)
 \le C_d\sum_i r_i^k
 \le C_dj\sum_i\mu(B_i)
 \le C_dj\,\varepsilon.
\]
先令覆盖尺度趋零，再令 \(\varepsilon\) 趋零，得到 \(\mathcal H^k(N_j)=0\)。这证明了\eqref{b1:eq:essential-reduced-boundary}。

最后处理 \(d=1\)。在任意内部有界区间 \(I=(a,b)\) 上，记 \(m=D\chi_E\)，并令 \(F(t)=m((a,t])\)。Fubini 定理直接给出 \(DF=m\)，因此 \(\chi_E-F\) 的分布导数为零，几乎处处等于常数。这样得到 \(\chi_E\) 的右连续有限变差代表 \(u\)。它处处只能取 \(0,1\)：从任一点的右侧选取避开原零测例外集的点列，再用右连续性即可。

每次从 \(0\) 变到 \(1\) 或反过来都消耗至少一个单位的变差，所以任一紧区间上只能有有限次这样的变化。于是 \(E\) 在零测集修改后局部为有限个区间的并，\(D\chi_E\) 在左端点取 \(+1\)、右端点取 \(-1\)，全变差就是这些端点的计数测度。它们恰为约化边界及测度论边界，密度为 \(1/2\)，外法向分别为 \(-1,+1\)。这给出所有结论的一维版本。
\end{proof}

证明的主干可与 Simon 的《Lectures on Geometric Measure Theory》定理14.3及其证明对读\cite[第14节，第72--76页]{Simon1983Lectures}。其中“约化点处的方向集中—缩放紧性—半空间极限”的顺序尤其重要。本书另外把图分片构造、密度识别及测度论边界余项的覆盖论证展开，以便这些结论只依赖已经建立的工具。

\subsection{周长测度的表示}
\label{b1:subsec:270504}

现在可以把前面的分布公式写成熟悉的面积通量形式。

\begin{theorem}[Gauss--Green]\label{b1:thm:gauss-green-finite-perimeter}
设 \(E\) 在 \(\Omega\) 中局部有限周长。对任意 \(\Phi\in C_c^1(\Omega;\mathbb R^d)\)，有
\begin{equation}\label{b1:eq:gauss-green-hausdorff}
 \int_E\operatorname{div}\Phi\,\mathrm dx
 =\int_{\partial^*E}\Phi\cdot\nu_E\,\mathrm d\mathcal H^{d-1}.
\end{equation}
对任意 Borel 集 \(A\subset\Omega\)，
\[
 P(E,A)=\mathcal H^{d-1}(A\cap\partial^*E)
       =\mathcal H^{d-1}(A\cap\partial^eE).
\]
\end{theorem}
\begin{proof}
将\eqref{b1:eq:de-giorgi-perimeter-measure}代入\eqref{b1:eq:gauss-green-perimeter-measure}，得到通量公式。周长的两种表示分别来自同一测度等式及\eqref{b1:eq:essential-reduced-boundary}。
\end{proof}

第二种周长表示不意味着 \(\partial^eE\) 的每一点都有法向。第一象限的角点已经给出反例。积分能够忽略这些点，是因为结构定理证明它们在 \(\mathcal H^{d-1}\) 下构成零集；不是在定义法向时可以随意省去任意坏点。

也不能把 \(\partial^*E\) 换成任意代表的普通边界。给一个光滑集合删去内部稠密零测集，不会改变 \(D\chi_E\)、周长、约化边界或测度论边界，却可以让普通边界包含整个原区域。因此有限周长理论的几何对象是由等价类恢复的边界，而不是某次点集表示附带的拓扑边界。

这一表示还说明，示性函数的导数是集中在余维一可整流集上的测度；一般 BV 函数的变化则由全部水平集共同组成。与\cref{b1:thm:bv-coarea}合起来，函数的总变差便可以理解为各层约化边界面积的累积。不过本节没有把所有 BV 导数都断言为单独一张超曲面上的面积：不同水平的边界可以形成完全不同的整体分布。

% !TeX root = ../../Book1.tex
\section{等周不等式}
\label{b1:sec:2706}

周长描述边界的代价，体积描述集合占据的空间。两者的联系不仅属于几何：
将一个函数拆成超水平集，集合的体积估计就能转化为函数的可积性估计。
反过来，把函数不等式用于示性函数，又会恢复集合不等式。
本节在全空间完整建立这两个方向；所证明的是只依赖维数的常数，并不预先宣称其最优。

Simon 的《Lectures on Geometric Measure Theory》\cite[\S6, pp. 34--40]{Simon1983Lectures}
把局部变差、磨光与 Poincaré 型估计放在同一条线上。
这里将前节的严格逼近接到\secref{b1:sec:2402}已经证明的端点 Sobolev 不等式，
再用余面积公式说明其集合意义。

% 实读范围：本地 Simon1983Lectures 正式扫描本，第6节印刷34--40页，
% 尤其6.1、6.2及6.4--6.6。该范围用于局部变差/磨光/Poincare语境，
% 不冒称提供最佳欧氏等周常数或球等号唯一性的证明。
% 主证明直接依赖本书24章GNS端点、27.2严格逼近、27.3余面积。
% 全空间Cc∞严格逼近在下面实际补做远处截断，不把任意域光滑逼近写成紧支集逼近。

\subsection{有界变差函数形式}
\label{b1:subsec:270601}

\begin{theorem}[有界变差函数的 Sobolev 端点估计]\label{b1:thm:bv-sobolev-endpoint}
设 \(d\ge2\)，\(q=d/(d-1)\)。存在仅依赖维数的常数 \(S_d<\infty\)，使
每个 \(u\in BV(\mathbb R^d)\) 都属于 \(L^q(\mathbb R^d)\)，且
\begin{equation}\label{b1:eq:bv-sobolev-endpoint}
 \|u\|_{L^q(\mathbb R^d)}\le S_d|Du|(\mathbb R^d).
\end{equation}
\end{theorem}
\begin{proof}
先把严格逼近改成全空间紧支集逼近。由定理%
\ref{b1:thm:bv-strict-approximation}，可取
\(w_j\in C^\infty(\mathbb R^d)\cap W^{1,1}(\mathbb R^d)\)，使
\[
 \|w_j-u\|_1\to0,\quad
 \int_{\mathbb R^d}|\nabla w_j|\,\mathrm dx
       \longrightarrow|Du|(\mathbb R^d).
\]
选光滑截断 \(0\le\chi_R\le1\)，在 \(B(0,R)\) 上等于 \(1\)，
支集包含于 \(B(0,2R)\)，且 \(|\nabla\chi_R|\le C/R\)。
对固定 \(j\)，当 \(R\to\infty\) 时，
\[
 \|\chi_Rw_j-w_j\|_1\to0,\quad
 \int_{\mathbb R^d}|w_j|\cdot|\nabla\chi_R|\,\mathrm dx
       \le\frac C R\|w_j\|_1\to0.
\]
因此可取 \(R_j\ge j\)，使这两个量都小于 \(1/j\)。
令 \(v_j=\chi_{R_j}w_j\)，则 \(v_j\in C_c^\infty(\mathbb R^d)\)，
\(v_j\to u\) 于 \(L^1\)，且由乘积公式，
\[
 \limsup_j\int_{\mathbb R^d}|\nabla v_j|\,\mathrm dx
 \le\lim_j\left(\int_{\mathbb R^d}|\nabla w_j|\,\mathrm dx+\frac1j\right)
 =|Du|(\mathbb R^d).
\]
下半连续性给出反向的下极限不等式，所以
\(\int|\nabla v_j|\to|Du|(\mathbb R^d)\)。

现在对 \(v_j\) 应用\cref{b1:thm:gns-endpoint}：
\[
 \|v_j\|_q\le S_d\int_{\mathbb R^d}|\nabla v_j|\,\mathrm dx.
\]
再从 \(L^1\) 收敛中抽取几乎处处收敛子列，利用 Fatou 引理，得到
\[
 \|u\|_q
 \le\liminf_j\|v_j\|_q
 \le S_d|Du|(\mathbb R^d).
\]
这一步同时证明目标 \(L^q\) 可积性，不要求预先知道 \(u\in L^q\)。
\end{proof}

右端没有零阶项，但左端函数类仍要求 \(u\in L^1(\mathbb R^d)\)。
这一条件使远处截断的误差消失，也排除了全空间上的非零常数。
严格逼近在这里提供的是总变差的数值收敛，不是导数测度的变差范数收敛；
证明也没有推出 \(v_j\to u\) 于临界空间 \(L^q\)。

维数 \(d=1\) 时，对 \(v\in C_c^\infty(\mathbb R)\) 分别从左右无穷远积分，可得
\[
 |v(x)|\le\int_{-\infty}^x|v'|,\quad
 |v(x)|\le\int_x^\infty|v'|,
 \quad
 2\|v\|_\infty\le\int_{\mathbb R}|v'|.
\]
沿上面同一组紧支集严格逼近，并在几乎处处极限上取界，得到
\[
 \|u\|_{L^\infty(\mathbb R)}\le\frac12|Du|(\mathbb R),
 \quad u\in BV(\mathbb R).
\]
这是单独的一维端点，不把 \(d/(d-1)\) 当作有限指数使用。

\subsection{有限周长集合形式}
\label{b1:subsec:270602}

\begin{theorem}[有限周长集合的等周估计]\label{b1:thm:finite-perimeter-isoperimetric}
设 \(d\ge2\)，集合 \(E\subseteq\mathbb R^d\) 可测、体积有限且周长有限。
则
\begin{equation}\label{b1:eq:finite-perimeter-isoperimetric}
 m_d(E)^{(d-1)/d}\le S_dP(E).
\end{equation}
\end{theorem}
\begin{proof}
有限体积给出 \(\mathbf1_E\in L^1(\mathbb R^d)\)，有限周长给出
\(\mathbf1_E\in BV(\mathbb R^d)\)。将它代入\eqref{b1:eq:bv-sobolev-endpoint}，
并注意
\[
 \|\mathbf1_E\|_{d/(d-1)}=m_d(E)^{(d-1)/d},
 \quad |D\mathbf1_E|(\mathbb R^d)=P(E),
\]
便得结论。
\end{proof}

这类将体积与边界大小联系起来的估计称为%
\term[dengzhou-budengshi]{等周不等式}[isoperimetric inequality]。
这里使用的是全空间周长。若将右端直接改为 \(P(E,\Omega)\)，
取 \(E=\Omega\) 就会出现正体积与零相对周长，因此必须另加适合区域的归一化。
有限体积也不可删除：全空间自身的周长为零，却有无穷体积。
在一维，前面的端点界说明 \(0<m_1(E)<\infty\) 的有限周长集合满足 \(P(E)\ge2\)；
一个有界非空区间达到这个值。

球可以用来检查常数的尺度。记 \(\omega_d=m_d(B(0,1))\)。
由光滑域周长公式\ref{b1:prop:smooth-domain-perimeter}及球面积计算，
\[
 m_d(B(a,r))=\omega_dr^d,\quad
 P(B(a,r))=d\omega_dr^{d-1}.
\]
所以任何适用于所有集合的常数 \(S_d\) 都必须满足
\[
 S_d\ge\frac1{d\omega_d^{1/d}}.
\]
半径恰好消去，说明指数 \((d-1)/d\) 与边界尺度相匹配。
但求出球的比值，并不证明这个下界对所有集合都可作为不等式常数，
更不证明哪些集合能取等号。本节不从非最优的 Gagliardo–Nirenberg–Sobolev 估计
推断最佳等周常数或等号分类。

\subsection{Sobolev 与等周不等式的联系}
\label{b1:subsec:270603}

从函数不等式到集合不等式只需代入示性函数。反方向则需要把函数的全部层重新组合，
余面积公式正好将累积的周长转回总变差。

\begin{proposition}[集合与函数估计的常数传递]\label{b1:prop:isoperimetric-sobolev-equivalence}
固定 \(d\ge2\)、\(q=d/(d-1)\) 及 \(C>0\)。下列两条等价。
\begin{equivalences}
 \item\label{b1:item:bv-isoperimetric-set-bound}
 每个有限体积的有限周长集合 \(E\subseteq\mathbb R^d\) 都满足
 \(m_d(E)^{1/q}\le CP(E)\)。
 \item\label{b1:item:bv-isoperimetric-function-bound}
 每个实值 \(u\in BV(\mathbb R^d)\) 都满足
 \(\|u\|_q\le C|Du|(\mathbb R^d)\)。
\end{equivalences}
\end{proposition}
\begin{proof}
由\itemref{b1:item:bv-isoperimetric-function-bound}推出%
\itemref{b1:item:bv-isoperimetric-set-bound}已在前一定理中证明。
现假设集合估计成立。

先说明可以取绝对值而不增加总变差。用严格逼近取
\(u_j\in W^{1,1}(\mathbb R^d)\)，使 \(u_j\to u\) 于 \(L^1\)，且
\(\int|\nabla u_j|\to|Du|(\mathbb R^d)\)。
Sobolev 绝对值公式给出 \(|u_j|\in W^{1,1}\)，
\(\bigl|\nabla|u_j|\bigr|\le|\nabla u_j|\) 几乎处处。
又有 \(\||u_j|-|u|\|_1\le\|u_j-u\|_1\)，故变差下半连续性推出
\[
 w=|u|\in BV(\mathbb R^d),\quad
 |Dw|(\mathbb R^d)\le|Du|(\mathbb R^d).
\]

令 \(E_t=\{w>t\}\)，\(t>0\)。由于 \(w\in L^1\)，有
\(m_d(E_t)\le\|w\|_1/t<\infty\)；余面积公式又保证几乎每个 \(E_t\) 周长有限。
对 \(R,N>0\)，用积分 Minkowski 不等式及集合估计，得到
\[
 \begin{aligned}
 \|\mathbf1_{B(0,R)}\min\{w,N\}\|_q
 &=\left\|\int_0^N\mathbf1_{E_t\cap B(0,R)}\,\mathrm dt\right\|_q\\
 &\le\int_0^N m_d\bigl(E_t\cap B(0,R)\bigr)^{1/q}\,\mathrm dt\\
 &\le\int_0^N m_d(E_t)^{1/q}\,\mathrm dt\\
 &\le C\int_0^\infty P(E_t)\,\mathrm dt
   =C|Dw|(\mathbb R^d).
 \end{aligned}
\]
空间与层值的截断保证最初的函数有界且支集有限，因而没有预设 \(w\in L^q\)。
令 \(R,N\to\infty\)，单调收敛定理给出
\[
 \|u\|_q=\|w\|_q\le C|Dw|(\mathbb R^d)\le C|Du|(\mathbb R^d).
\]
同一个常数 \(C\) 在两个方向均未改变，证明完成。
\end{proof}

这个等价性说明，最佳集合常数与最佳有界变差函数常数若按各自下确界定义，必然相同；
但它本身并不计算那个下确界。
在本书当前的证明链中，先由坐标积分建立 Sobolev 估计，再由严格逼近与余面积转到集合，
所有实际使用的常数已经得到控制。若转而从尖锐的几何等周定理出发，
同一层积分论证也能将它的常数原样传给有界变差函数的端点估计。


\begin{chaptersummary}
有界变差函数以有限导数测度控制变化，既容纳可积梯度，也容纳跳跃及更分散的奇异变化。
光滑逼近保持函数的 \(L^1\) 极限和总变差数值，却通常不能保证导数测度在变差范数中收敛。
紧性来自平移控制。本章在任意开集上证明的全域版本还要求质量不逃向边界或无穷远。
对有界 Lipschitz 区域，标准紧嵌入可由另需建立的有界变差函数延拓定理得到，不必另加这种尾部假设。

余面积公式把导数变差分解为水平集周长。对每个有限周长集合，
De Giorgi 定理进一步把周长识别成约化边界的余维一面积，
其证明先取得非退化半空间极限，再构造可数 Lipschitz 图并辨认密度。
Gauss–Green 公式中的外法向负号须与分布导数的约定配合；换用其他文献的记号时，也应同步调整。

全空间有界变差函数的端点不等式与有限体积集合的等周不等式具有相同的常数传递机制。
本章给出有效常数，不讨论最佳常数。
一般有界变差函数并非单个集合的示性函数；各层边界的累积还会形成怎样的导数，
将由下一章的精细结构解释。
\end{chaptersummary}

\BookExercises

% \chapref{b1:ch:27}习题临时稿；不另设 BookExercises 入口。
% 九题均由本章与前章已证明工具自拟，逐题注明实质增量，不冒认教材原题号。
% 不调用第28章的一般有界变差函数链式法则、跳跃集结构或特殊有界变差函数理论。

\begin{exercise}\label{b1:ex:27-grid-interface-cancellation}
取正整数 \(m,n\)，令 \(\Omega=(0,m)\times(0,n)\)。
在每个单位格
\[
 Q_{ij}=(i-1,i)\times(j-1,j),\quad
 1\le i\le m,\quad1\le j\le n,
\]
上令 \(u=a_{ij}\)，其中 \(a_{ij}\in\mathbb R\)。
网格线上的值任取。求 \(Du\) 和 \(|Du|(\Omega)\)，并求全空间零延拓 \(E_0u\) 的总变差。
解释为什么内部界面的贡献是相邻值之差，而不是两个值的绝对值之和。
\end{exercise}
% 来源：自拟；把单个矩形的导数推广到共享界面，重点是先相消再取向量变差。
\begin{solution}
对光滑紧支集测试函数逐格积分。固定一条竖界面
\(V_{ij}=\{i\}\times(j-1,j)\)，左格在右侧贡献
\(-a_{ij}e_1\mathcal H^1\)，右格在左侧贡献
\(a_{i+1,j}e_1\mathcal H^1\)。
水平界面 \(H_{ij}=(i-1,i)\times\{j\}\) 同理。因此在 \(\Omega\) 中，
\[
 \begin{aligned}
 Du={}&\sum_{i=1}^{m-1}\sum_{j=1}^n
       (a_{i+1,j}-a_{ij})e_1\,\mathcal H^1\!\llcorner V_{ij}\\
 &+\sum_{i=1}^m\sum_{j=1}^{n-1}
       (a_{i,j+1}-a_{ij})e_2\,\mathcal H^1\!\llcorner H_{ij}.
 \end{aligned}
\]
不同线段只可能在端点相交，这些点的 \(\mathcal H^1\) 测度为零。
每条线段长度为一，所以
\[
 |Du|(\Omega)=
 \sum_{i=1}^{m-1}\sum_{j=1}^n|a_{i+1,j}-a_{ij}|
 +\sum_{i=1}^m\sum_{j=1}^{n-1}|a_{i,j+1}-a_{ij}|.
\]
外边界在相对导数中没有贡献。补零后，下侧与左侧的系数为内部值，
上侧与右侧的系数为内部值的负号；它们与内部界面互不重叠到面测度零集。
于是
\[
 \begin{aligned}
 |D(E_0u)|(\mathbb R^2)
 ={}&|Du|(\Omega)
    +\sum_{j=1}^n(|a_{1j}|+|a_{mj}|)\\
   &+\sum_{i=1}^m(|a_{i1}|+|a_{in}|).
 \end{aligned}
\]
若相邻两格取同一个值，跨界面并没有变化，故该界面必须完全抵消。
先对每格分别取变差再相加，会丢失导数测度的这个抵消过程。
\end{solution}

\begin{exercise}\label{b1:ex:27-small-balls-perimeter-threshold}
设 \(d\ge2\)、\(\alpha>1/d\)，并令
\[
 x_j=3je_1,\quad r_j=\frac14j^{-\alpha},\quad
 E=\bigcup_{j=1}^\infty B(x_j,r_j).
\]
证明 \(m_d(E)<\infty\)，并确定哪些 \(\alpha\) 使 \(E\) 有有限周长。
在允许无穷的意义下计算 \(P(E)\)，说明有限体积不能控制周长。
\end{exercise}
% 来源：自拟；用局部有限且相隔的小球把体积和边界的两个尺度化成不同级数。
\begin{solution}
这些闭球两两不交，而且每个紧集只与有限个球相交。
由体积的可数可加性，
\[
 m_d(E)=\omega_d4^{-d}\sum_{j=1}^\infty j^{-\alpha d}<\infty.
\]
对任意紧支集测试场，只有有限个球参与积分。
光滑域周长公式因而给出局部向量测度
\[
 D\mathbf1_E
 =-\sum_{j=1}^\infty
      \nu_j\,\mathcal H^{d-1}\!\llcorner\partial B(x_j,r_j),
 \quad \nu_j(x)=\frac{x-x_j}{r_j}.
\]
这是局部有限的和；各球面互不相交，故变差没有相消，且
\[
 P(E)=d\omega_d\sum_{j=1}^\infty r_j^{d-1}
     =d\omega_d4^{1-d}\sum_{j=1}^\infty j^{-\alpha(d-1)}.
\]
若级数收敛，右侧给出全空间有限导数测度，所以 \(E\) 有有限周长。
若级数发散，包含前 \(N\) 个球面的有界开集上的变差至少为前 \(N\) 项之和，
让 \(N\to\infty\) 即知全空间周长无穷。
因此
\[
 P(E)<\infty\quad\Longleftrightarrow\quad
       \alpha>\frac1{d-1}.
\]
特别地，\(1/d<\alpha\le1/(d-1)\) 时体积有限而周长无限。
两个端点不同，是因为球体积按 \(r^d\) 缩放，而球周长按 \(r^{d-1}\) 缩放。
\end{solution}

\begin{exercise}\label{b1:ex:27-basic-boundaries-and-blowup}
分别考察下列集合：
\[
 H=\{x\in\mathbb R^d:x_d<0\},\qquad
 Q=(0,\infty)^2\subset\mathbb R^2,
\]
以及
\[
 F=B(0,1)\setminus\mathbb Q^d\subset\mathbb R^d.
\]
对每个集合求其密度零点集、密度一点集、测度论边界与约化边界，并在约化边界上求测度论外法向。比较 \(F\) 的拓扑边界与测度论边界。最后直接验证：半空间在边界点的缩放恒为相应半空间；球在每个边界点的集合缩放与周长测度缩放分别趋于切半空间与切平面测度。说明这些结论为何也适用于 \(F\)。
\end{exercise}
% 来源：自拟；用半空间、角点及稠密零测修改训练三种边界与缩放极限的基本辨析。
\begin{solution}
先看半空间。直接由球的对称性，
\[
 H^1=H,\qquad H^0=\{x_d>0\},\qquad
 \partial^eH=\partial^*H=\{x_d=0\}.
\]
在边界上，外法向为 \(\nu_H=e_d\)，并且
\[
 D\mathbf1_H=-e_d\,\mathcal H^{d-1}
                    \!\llcorner\{x_d=0\}.
\]
若 \(x\in\partial H\)，则 \(H_{x,r}=H\) 对每个 \(r>0\) 都成立；相应的缩放周长测度也始终就是 \(\mathcal H^{d-1}\llcorner\{x_d=0\}\)。

对第一象限，
\[
 Q^1=Q,\qquad Q^0=\mathbb R^2\setminus\overline Q,\qquad
 \partial^eQ=\bigl([0,\infty)\times\{0\}\bigr)
              \cup\bigl(\{0\}\times[0,\infty)\bigr).
\]
正半轴上的非原点处密度为 \(1/2\)，原点处密度为 \(1/4\)，所以它们都属于测度论边界。两条开正半轴上的外法向分别为
\[
 \nu_Q(x_1,0)=-e_2\quad(x_1>0),qquad
 \nu_Q(0,x_2)=-e_1\quad(x_2>0).
\]
在原点，两个边界射线对 \(D\mathbf1_Q\) 的贡献分别沿 \(e_2\) 与 \(e_1\)，故
\[
 \frac{D\mathbf1_Q(B(0,r))}{|D\mathbf1_Q|(B(0,r))}
 =\frac{e_1+e_2}{2},
\]
其长度为 \(1/\sqrt2\)，不是 \(1\)。因此
\[
 \partial^*Q=\bigl((0,\infty)\times\{0\}\bigr)
              \cup\bigl(\{0\}\times(0,\infty)\bigr),
\]
角点属于 \(\partial^eQ\setminus\partial^*Q\)。

集合 \(F\) 与单位球只相差可数零测集，所以它们的密度点、分布导数、周长测度和约化边界完全相同：
\[
 F^1=B(0,1),\qquad F^0=\mathbb R^d\setminus\overline{B(0,1)},\qquad
 \partial^eF=\partial^*F=\partial B(0,1),\qquad
 \nu_F(x)=x.
\]
另一方面，\(F\) 与其补集都在单位球内稠密，故 \(F\) 的内部为空、闭包为 \(\overline{B(0,1)}\)，从而它的拓扑边界是整个闭球。这说明拓扑边界不具有零测修改不变性。

最后，对 \(x\in\partial B(0,1)\)，有
\[
 z\in B(0,1)_{x,r}
 \quad\Longleftrightarrow\quad
 z\cdot x<-\frac r2|z|^2.
\]
在每个固定球内，右侧阈值一致趋于零，边界超平面又是零测集，故示性函数在局部 \(L^1\) 中趋于 \(\mathbf1_{\{z\cdot x<0\}}\)。旋转到 \(x=e_d\) 后，靠近原点的缩放球面写成图
\[
 z_d=\frac{-1+\sqrt{1-r^2|z'|^2}}{r},
\]
它及其图面积因子在紧集上一致趋于 \(0\) 与 \(1\)；球面的另一支在固定紧集之外。由图面积公式或直接参数换元，缩放周长测度趋于
\(\mathcal H^{d-1}\llcorner x^\perp\)。由于 \(F\) 与球有相同的示性函数等价类及周长测度，同样的两种缩放结论也适用于 \(F\)。
\end{solution}

\begin{optionalexercise}\label{b1:ex:27-translation-bv-norm-characterization}
设 \(u\in BV(\mathbb R^d)\)，\(\tau_hu(x)=u(x-h)\)。
\begin{enumerate}
 \item\label{b1:item:27-translation-strict}
 证明当 \(h\to0\) 时，\(\tau_hu\) 总是严格收敛到 \(u\)。
 \item\label{b1:item:27-translation-norm}
 证明
 \[
  \|\tau_hu-u\|_{BV(\mathbb R^d)}\longrightarrow0
    \quad\Longleftrightarrow\quad u\in W^{1,1}(\mathbb R^d).
 \]
 \item\label{b1:item:27-translation-interval}
 对 \(u=\mathbf1_{(0,1)}\)，计算 \(0<h<1\) 时的
 \(\|\tau_hu-u\|_{BV(\mathbb R)}\)。
\end{enumerate}
\end{optionalexercise}
% 来源：自拟；以测度在全变差范数中的平移连续性辨认绝对连续导数，
% 不重复正文对阶跃函数作光滑严格逼近的例子。
\begin{solution}
函数的 \(L^1\) 平移连续，而
\[
 D(\tau_hu)=\tau_hDu,\quad
 |D(\tau_hu)|(\mathbb R^d)=|Du|(\mathbb R^d),
\]
其中 \((\tau_h\mu)(A)=\mu(A-h)\)。
这证明第一项。若 \(u\in W^{1,1}\)，则函数及各弱偏导数都在 \(L^1\) 中平移连续，
从而得到第二项的反向蕴涵。

现假设平移在 \(BV\) 范数中连续，记 \(\mu=Du\)。
于是 \(|\tau_h\mu-\mu|(\mathbb R^d)\to0\)。
用非负单位积分核磨光测度，记
\(\mu_\varepsilon=(\rho_\varepsilon*\mu)m_d\)。
对每个连续紧支集向量场 \(\Phi\)，Fubini 定理给出
\[
 \int\Phi\cdot\mathrm d(\mu_\varepsilon-\mu)
 =\int\rho_\varepsilon(h)
       \left(\int\Phi\cdot\mathrm d(\tau_h\mu-\mu)\right)\,\mathrm dh.
\]
向量测度全变差的连续测试对偶式于是给出
\[
 |\mu_\varepsilon-\mu|(\mathbb R^d)
 \le\sup_{|h|\le\varepsilon}
          |\tau_h\mu-\mu|(\mathbb R^d)\longrightarrow0.
\]
每个 \(\mu_\varepsilon\) 都对 Lebesgue 测度绝对连续。
若 \(N\) 为 Lebesgue 零测 Borel 集，则
\[
 |\mu|(N)\le|\mu-\mu_\varepsilon|(\mathbb R^d)
                     +|\mu_\varepsilon|(N)\longrightarrow0.
\]
所以 \(\mu\ll m_d\)。各分量由 Radon–Nikodym 定理写为 \(L^1\) 密度，
这些密度正是 \(u\) 的弱偏导数，因而 \(u\in W^{1,1}\)。
此处连续测试对偶式由极分解和连续紧支集函数的积分逼近得到，与变差对偶证明中的步骤相同。

最后，当 \(u=\mathbf1_{(0,1)}\) 时，
\[
 \|\tau_hu-u\|_1=2h,\quad
 D(\tau_hu-u)=\delta_h-\delta_{1+h}-\delta_0+\delta_1.
\]
四个点互不相同，故导数差的全变差为 \(4\)，于是
\[
 \|\tau_hu-u\|_{BV(\mathbb R)}=2h+4.
\]
函数与自身的小平移严格接近，却没有在 \(BV\) 范数中接近。
\end{solution}

\begin{optionalexercise}\label{b1:ex:27-finite-valued-strict-approximation}
设 \(m_d(\Omega)<\infty\)，\(u\in BV(\Omega)\)，且 \(0\le u\le M<\infty\) 几乎处处。
证明存在只取有限多个值的 \(v_j\in BV(\Omega)\)，使
\[
 \|v_j-u\|_{L^\infty(\Omega)}\longrightarrow0,\quad
 |Dv_j|(\Omega)\longrightarrow|Du|(\Omega).
\]
这里 \(L^\infty\) 指本性上确界范数。要求将每个 \(v_j\) 写成有限周长超水平集的示性函数之和，
并说明层值怎样选择。
\end{optionalexercise}
% 来源：自拟；在一组平移的等距层值中平均周长，选出变差受控的量化函数。
\begin{solution}
固定 \(h>0\)，对 \(0<\theta<h\) 定义
\[
 v_{h,\theta}
   =h\sum_{k=0}^\infty\mathbf1_{\{u>\theta+kh\}}.
\]
因为 \(u\le M\)，非空层只有有限个，所得函数取有限多个值。
对每个 \(s\in[0,M]\)，等距网格的计数误差小于或等于一个网格长度，
所以
\[
 |v_{h,\theta}-u|\le h\quad\text{几乎处处}.
\]
余面积公式说明，对几乎每个 \(\theta\)，上式涉及的全部层集都有有限周长；
只需排除每个平移网格命中坏层值的零测集，再取可数并。
在这些 \(\theta\) 上，
\[
 |Dv_{h,\theta}|(\Omega)
 \le S_h(\theta)=
        h\sum_{k=0}^\infty P(\{u>\theta+kh\},\Omega).
\]
由 Tonelli 定理及 \(u\ge0\)，
\[
 \frac1h\int_0^hS_h(\theta)\,\mathrm d\theta
 =\int_0^\infty P(\{u>t\},\Omega)\,\mathrm dt
 =|Du|(\Omega).
\]
因此可选一个容许的 \(\theta_h\)，使
\(S_h(\theta_h)\le|Du|(\Omega)+h\)。
取 \(h_j\downarrow0\)，令 \(v_j=v_{h_j,\theta_{h_j}}\)。
有限体积保证
\[
 \|v_j-u\|_1\le h_jm_d(\Omega)\longrightarrow0,
 \quad
 \limsup_j|Dv_j|(\Omega)\le|Du|(\Omega).
\]
下半连续性给出相反的下极限不等式，因而总变差收敛。
有限值逼近与光滑逼近采用不同的对象：前者以有限周长层集承载界面，
后者以光滑梯度承载变化；二者都只要求总变差数值收敛。
\end{solution}

\begin{exercise}\label{b1:ex:27-irregular-zero-extension}
令
\[
 a_j=2^{-j},\quad b_j=2^{-j}+2^{-j-2},\quad
 \Omega=\bigcup_{j=1}^\infty(a_j,b_j)\subset(0,1).
\]
在 \(\Omega\) 上取 \(u=1\)，并令 \(u_N\) 在前 \(N\) 个区间上等于 \(1\)，
在其余区间上等于 \(0\)。
证明 \(u_N\to u\) 于 \(BV(\Omega)\)，每个全空间零延拓 \(E_0u_N\) 都属于
\(BV(\mathbb R)\)，但 \(E_0u\notin BV(\mathbb R)\)。
计算 \(|D(E_0u_N)|(\mathbb R)\)。
\end{exercise}
% 来源：自拟；域内没有导数的常数分片，经补零后产生无穷多个端点代价。
\begin{solution}
这些区间两两分离，且长度为 \(2^{-j-2}\)。
函数 \(u_N\) 和 \(u\) 在每个连通分量上为常数，所以它们在 \(\Omega\) 中的分布导数均为零。
因此
\[
 \|u_N-u\|_{BV(\Omega)}
 =\sum_{j>N}(b_j-a_j)\longrightarrow0.
\]
另一方面，在全空间中逐区间计算，得到
\[
 D(E_0u_N)=\sum_{j=1}^N(\delta_{a_j}-\delta_{b_j}),
 \quad
 |D(E_0u_N)|(\mathbb R)=2N.
\]
所有端点互不相同，因此没有相消。

若 \(E_0u\) 有有限变差，则在每个端点的小邻域内，它的导数都必须与上述单个区间的导数一致。
每次固定有限多个端点，可以选取两两不交的小邻域并避开其余端点。
故对每个 \(N\)，其总变差至少为 \(2N\)，矛盾。
这个例子不否认内部紧支集截断的零延拓性质：\(u\) 的非零部分一直延伸到
无穷多个分量端点，没有同一个内部紧集隔开这些边界。
\end{solution}

\begin{exercise}\label{b1:ex:27-eccentric-hole-flux}
在 \(\mathbb R^2\) 中设 \(R,r>0\)、\(|a|+r<R\)，并令
\[
 E=B(0,R)\setminus\overline{B(a,r)}.
\]
求 \(E\) 的外法向、\(D\mathbf1_E\) 及 \(P(E)\)。
再对向量场
\[
 F(x)=\frac{x-a}{|x-a|^2}
\]
计算两条边界曲线各自的外向通量，解释为什么总通量为零。
不要求使用绕数或复分析。
\end{exercise}
% 来源：自拟；用偏心孔检验内边界方向，以光滑域Gauss计算奇点被排除后的通量。
\begin{solution}
外圆上的外法向为 \(x/R\)，内圆上的外法向指向孔内，即
\[
 \nu_E(x)=
 \begin{cases}
 x/R,&x\in\partial B(0,R),\\
 -(x-a)/r,&x\in\partial B(a,r).
 \end{cases}
\]
由光滑域周长公式，
\[
 D\mathbf1_E
 =-\frac{x}{R}\,\mathcal H^1\!\llcorner\partial B(0,R)
   +\frac{x-a}{r}\,\mathcal H^1\!\llcorner\partial B(a,r),
 \quad P(E)=2\pi(R+r).
\]
两条边界的周长相加，虽然它们的外法向相对于各自圆心取了不同的方向。

场 \(F\) 在 \(\overline E\) 的邻域内光滑，直接求导得
\[
 \operatorname{div}F(x)
 =\frac2{|x-a|^2}-\frac{2|x-a|^2}{|x-a|^4}=0.
\]
如需紧支集测试场，可乘以在 \(\overline E\) 的邻域等于 \(1\)、支集避开 \(a\) 的光滑截断。
在内圆上，\(F\cdot\nu_E=-1/r\)，故内圆通量为 \(-2\pi\)。
Gauss 公式给出两条曲线的通量之和等于 \(\int_E\operatorname{div}F=0\)，
所以外圆通量为 \(2\pi\)。
奇点位于被删去的孔中，不在积分区域内；把内孔也按远离其圆心的方向取法向，
会错误地丢失这个抵消。
\end{solution}

\begin{exercise}\label{b1:ex:27-signed-measure-primitive}
设 \(-\infty<a<b<\infty\)，\(\mu\) 是 \((a,b)\) 上的有限 Radon 符号测度，
并令
\[
 F(x)=\mu\bigl((a,x]\bigr),\quad a<x<b,\quad M=\mu\bigl((a,b)\bigr).
\]
\begin{enumerate}
 \item\label{b1:item:27-measure-primitive}
 证明 \(F\in BV(a,b)\)、\(DF=\mu\)，并证明所有分布导数为 \(\mu\) 的
 有界变差函数都与某个 \(F+c\) 几乎处处相等。
 \item\label{b1:item:27-measure-primitive-endpoints}
 求 \(F\) 的左右极限及跳跃值。将 \(F+c\) 在 \((a,b)\) 外补零，记为 \(G_c\)，
 计算 \(DG_c\) 和 \(|DG_c|(\mathbb R)\)。
 \item\label{b1:item:27-measure-primitive-optimal-constant}
 在全部常数 \(c\in\mathbb R\) 中，求上述零延拓总变差的最小值。
\end{enumerate}
\end{exercise}
% 来源：自拟；一般符号测度原函数及端点代价，不复述正文的Cantor概率测度例子。
\begin{solution}
有 \(|F(x)|\le|\mu|\bigl((a,b)\bigr)\)，故 \(F\in L^1(a,b)\)。
对 \(\varphi\in C_c^\infty(a,b)\)，Fubini 定理给出
\[
 \int_a^bF(x)\varphi'(x)\,\mathrm dx
 =\int_{(a,b)}\left(\int_t^b\varphi'(x)\,\mathrm dx\right)\,\mathrm d\mu(t)
 =-\int_{(a,b)}\varphi\,\mathrm d\mu.
\]
换序由有限变差及 \((b-a)\|\varphi'\|_\infty\) 的控制保证。
因此 \(DF=\mu\)，且 \(|DF|\bigl((a,b)\bigr)=|\mu|\bigl((a,b)\bigr)\)。

若 \(Dv=\mu\)，令 \(w=v-F\)，则 \(Dw=0\)。
任何积分为零的 \(\psi\in C_c^\infty(a,b)\) 都可写成某个
\(\phi\in C_c^\infty(a,b)\) 的导数，故 \(\int w\psi=0\)。
固定 \(\eta\in C_c^\infty(a,b)\)、\(\int\eta=1\)，令 \(c=\int w\eta\)。
将任意测试函数减去其积分乘以 \(\eta\)，可得
\(\int w\psi=c\int\psi\)，于是正则分布的唯一性给出 \(w=c\) 几乎处处。

有限测度的集合连续性给出
\[
 F(a+)=0,\quad F(b-)=M,\quad
 F(x+)=F(x),\quad F(x-)=\mu\bigl((a,x)\bigr).
\]
所以内部跳跃值为 \(F(x)-F(x-)=\mu(\{x\})\)。
对全空间测试函数重复上面的换序，但保留端点项，得到
\[
 DG_c=\mu+c\delta_a-(c+M)\delta_b.
\]
其中 \(\mu\) 在区间外补零。内部测度与两个端点质量相互奇异，故
\[
 |DG_c|(\mathbb R)=|\mu|\bigl((a,b)\bigr)+|c|+|c+M|.
\]
由三角不等式，\(|c|+|c+M|\ge|M|\)，
并且当且仅当 \(c\) 位于连接 \(0\) 与 \(-M\) 的闭区间内时取等号。
故最小值为
\[
 |\mu|\bigl((a,b)\bigr)+|M|.
\]
内部总变差由导数固定，常数的选择只改变两个端点之间分配的补零代价。
\end{solution}

\begin{optionalexercise}\label{b1:ex:27-confined-bv-minimization}
设 \(d\ge2\)，考虑
\[
 \mathcal A=\left\{\,u\in BV(\mathbb R^d):
      u\ge0\ \text{几乎处处},\quad\int_{\mathbb R^d}u\,\mathrm dx=1\,\right\},
\]
以及允许取值 \(+\infty\) 的泛函
\[
 J(u)=|Du|(\mathbb R^d)+\int_{\mathbb R^d}|x|u(x)\,\mathrm dx.
\]
证明 \(J\) 在 \(\mathcal A\) 上取得有限的极小值。
再证明，若删除最后的一阶矩项，只最小化总变差，则下确界为零但不能取得。
指出矩项在极限过程中保持单位质量的具体作用。
\end{optionalexercise}
% 来源：自拟；全空间BV直接法由矩控制防止质量逃逸，不使用最佳等周常数或解的正则性。
\begin{solution}
取非负光滑紧支集函数并归一化积分，得到 \(J<\infty\) 的容许函数。
故下确界 \(I\) 有限且非负。选 \(u_j\in\mathcal A\)，使 \(J(u_j)\to I\)，
并可要求 \(J(u_j)\le I+1\)。
于是 \(\|u_j\|_1=1\)、总变差一致有界，且
\[
 \int_{|x|>R}u_j(x)\,\mathrm dx
 \le R^{-1}\int|x|u_j(x)\,\mathrm dx\le\frac{I+1}{R}.
\]
由带尾部控制的有界变差函数紧性%
\cref{b1:cor:bv-global-compactness-with-tails}，可取全空间 \(L^1\) 收敛子列，
记极限为 \(u\)。再抽几乎处处收敛子列，有 \(u\ge0\)，而全域 \(L^1\) 收敛保证
\(\int u=1\)，并非只得到 \(\int u\le1\)。
变差下半连续性和 Fatou 引理分别给出
\[
 |Du|(\mathbb R^d)\le\liminf_j|Du_j|(\mathbb R^d),\quad
 \int|x|u(x)\,\mathrm dx\le\liminf_j\int|x|u_j(x)\,\mathrm dx.
\]
所以 \(u\in\mathcal A\)，且 \(J(u)\le\liminf_jJ(u_j)=I\)，即取得极小值。

删除矩项后，取固定的非负 \(\eta\in C_c^\infty(\mathbb R^d)\)、\(\int\eta=1\)，令
\[
 v_R(x)=R^{-d}\eta(x/R).
\]
则 \(v_R\in\mathcal A\)，且
\[
 |Dv_R|(\mathbb R^d)=R^{-1}\|\nabla\eta\|_1\longrightarrow0.
\]
因此总变差的下确界为零。如果某个 \(v\in\mathcal A\) 达到它，
由\hyperref[b1:thm:bv-sobolev-endpoint]{有界变差函数的 Sobolev 端点估计\ref*{b1:thm:bv-sobolev-endpoint}\CJKecglue}便有 \(\|v\|_{d/(d-1)}=0\)，
与 \(\int v=1\) 矛盾。
膨胀序列在每个有界区域内的质量趋零，但总质量始终为一；
矩项所给的统一尾部估计，恰好排除了这种单位质量在局部极限中丢失的情形。
\end{solution}
